Enterprise Modelling and Information Systems Architectures
Vol. 17, No. 2 (2022). DOI:10.18417/emisa.17.2
Combining Goal modelling with Business Process modelling 1
Special Issue on 100 Years of Graphical Business Process Modelling by R. Laue, H. Mayr & B. Thalheim
Combining Goal modelling with Business Process modelling
Two Decades of Experience with the User Requirements Notation
Standard
Daniel Amyot*,a,b, Okhaide Akhigbea
, Malak Baslymanc
, Sepideh Ghanavatid
,
Mahdi Ghasemia
, Jameleddine Hassinec
, Lysanne Lessarda,b, Gunter
Mussbachere
, Kai Shena
, Eric Yuf
a University of Ottawa, Ottawa, Canada.
b
Institut du savoir Montfort, Ottawa, Canada.
c King Fahd University of Petroleum and Minerals, Dhahran, Saudi Arabia.
d University of Maine, Orono, USA.
e McGill University, Montreal, Canada
f University of Toronto, Toronto, Canada
Abstract. Goal modelling aims to capture stakeholder and system goals, together with social, intentional,
and structural relationships, in a way that supports trade-off analysis and decision making. Goal models
and business process models provide complementary and synergetic views of a system, which lead to
a more complete understanding of what exists and a better description of what needs to be designed
than with only one of these views. The User Requirements Notation (URN), standardized in 2008 by the
International Telecommunication Union (ITU-T) with improved versions in 2012 and 2018, combines goal
modelling with process modelling, while providing graphical and textual syntaxes for both views. URN
helps modellers exploit the value of social modelling in the context of process design and improvement. In
this paper, we report on nearly two decades of combined goal/process modelling with URN in different
areas – based on coarse-grained statistics from a literature review and the authors’ personal experiences
with URN. In particular, beyond the telecommunication services for which URN was created, we highlight
applications to goal/process alignments, regulatory compliance and intelligence, process adaptation and
improvement, value co-creation and service systems, and goal-oriented process mining, as well as several
advanced modelling techniques. An overview of the tool-supported analyses (for satisfaction, alignment,
compliance, and others) and model creation mechanisms (composition, aspects, slicing, adaptation, reuse,
and others) is also provided. The last part of this paper focuses on important challenges and exciting
opportunities for future research, especially in the areas of data-driven applications (e.g, AI/machine
learning), socio-cyber-physical systems, and usable automation.
Keywords. Goal-oriented modelling • GRL • Process modelling • UCM • User Requirements Notation
Communicated by Ralf Laue. Received 2021-03-15. Accepted after 1 revision on 2021-11-05.
* Corresponding author.
E-mail. damyot@uottawa.ca
This paper is dedicated to Raymond Buhr, inventor of the
Use Case Maps notation, who passed away in January 2021.
We thank the many people involved in the development of
URN and of related methods, standards, and tools since its
1 Introduction
Process modelling languages have now been
around for a century, and the seminal work on Proinception, as well as the anonymous referees for their useful
feedback. This work was supported in part by NSERC.
International Journal of Conceptual Modeling
Vol. 17, No. 2 (2022). DOI:10.18417/emisa.17.2
2 D. Amyot, O. Akhigbe, M. Baslyman, S. Ghanavati, M. Ghasemi, J. Hassine, L. Lessard, et al.
Special Issue on 100 Years of Graphical Business Process Modelling by R. Laue, H. Mayr & B. Thalheim
cess Charts by Gilbreth and Gilbreth (1921) has
since influenced dozens of variants and successors
(Mili et al. 2010). These domain-specific modelling languages (Karagiannis et al. 2016) enable
the representation and analysis of the activities of a
process (what), their sequencing (when), their subactivities (how), and their performers (who/where).
While contemporary process modelling languages
such as the Business Process Model and Notation
(BPMN) (OMG 2014), Activity Diagrams (OMG
2017), Petri Nets (ISO 2019), and Integration
Definition 0 (IDEF0) (IEEE 2012) can help engineers and business analysts reason about what,
when, where, who, and how questions, they are
rather limited in addressing why questions that are
important for trade-off analysis and for rationale
documentation, among others.
Goal modelling languages, on the other hand,
excel at describing rationales of systems and processes, together with social aspects and dependencies. These languages typically describe goals
(what/why), their refinements into activities (how),
and various relationships to performers (who and
for whom) and to each others (why). They thus provide important information about the who’s and
why’s of a system or organizational context, but
they do not express sequencing, which is one other
major difference they have with process modelling
languages.
Over the past 25 years, different goal modelling languages have been proposed (Horkoff et al.
2019), including i* (Dalpiaz et al. 2016; Yu 1997;
E. S. Yu et al. 2011), the Non-Functional Requirement Framework (Chung et al. 2000), GBRAM
(Antón 1996), KAOS (van Lamsweerde 2001,
2009), the Business Intelligence Model (BIM)
(Horkoff et al. 2014), and Tropos (Bresciani
et al. 2004). Some of these languages also have
their own extensions and variants. For example,
Gonçalves et al. (2018) have reviewed 96 papers
presenting extensions of the i* language. Many
of these languages have industrial applications as
well (Horkoff et al. 2019; E. Yu et al. 2013).
1.1 Process/Goal Combinations
Given that process modelling languages and
goal modelling languages have complementary
strengths regarding their modelling concepts and
the types of analysis questions their models can
answer, it becomes important to explore whether
combining their capabilities can be beneficial.
Over the years, several combinations have been
explored, including this non-exhaustive list:
• Three decades ago, Godart et al. (1992) explored early ideas about combining goal modelling with an informal process modelling notation for specifying cooperative software development activities. Goals were described as
logical predicates with positive and negative dependencies on activities, and a reasoning engine
enabled simple planning.
• Jacobs and Holten (1995) provided one of the
first integrations of simple process models with
goal models based on a common metamodel
implemented using Telos (Mylopoulos et al.
1990). This combination enabled goals to drive
decisions among explicit alternative activities
at the process level. Contributions between
goals were based on The House of Quality
(Clausing and Hauser 1988), with support for
fuzzy qualitative reasoning.
• Decreus et al. (2009) published a literature
review of practical challenges in transforming
goal models (in i* here) to business process
models in general (including BPMN). Many of
these challenges, related to concept mapping,
structures, sequencing, and automation, are still
current a decade later.
• Transformations in the opposite direction were
also explored. For example, Odeh et al. (2018)
investigated how to derive i* models (manually)
from a category of business process models
that are characterized as goal-based and roleoriented. How to do this for general process
models remains an issue.
• In a healthcare context, Cardoso et al. (2011) explored traceability and harmonization between
Tropos goal models and ARIS-based process
Enterprise Modelling and Information Systems Architectures
Vol. 17, No. 2 (2022). DOI:10.18417/emisa.17.2
Combining Goal modelling with Business Process modelling 3
Special Issue on 100 Years of Graphical Business Process Modelling by R. Laue, H. Mayr & B. Thalheim
models in order to improve goal-process alignment.
• Such alignment was also investigated by Koliadis and Ghose (2006), who have proposed
a method that combines KAOS goal models
with BPMN process models, with alignment
verification in changing, dynamic contexts.
• Gol Mohammadi (2019) loosely combined
BPMN with i* to model trustworthiness requirements in cyber-physical systems, with explicit
concepts for trustworthiness goals, concerns,
and threats.
These examples show that different combinations were explored to serve specific purposes or
domains. However, most approaches faced many
challenges related to the combination of languages
and tools that were not meant to interoperate in
the first place. This, in turn, limits their analysis
capabilities and potential benefits.
1.2 User Requirements Notation
One particular combination of two languages,
namely Use Case Maps (UCM) for process modelling and the Goal-oriented Requirement Language (GRL) for goal modelling, resulted in the
User Requirements Notation (URN). URN is a
standard from the International Telecommunication Union (ITU-T) that explicitly addresses goals
and processes in graphical and textual ways, in one
unified conceptual modelling language, with links
between the two perspectives (ITU-T 2018). This
standard also includes original features such as
measurable indicators and quantitative/qualitative
analysis for goal models, as well as executable
scenario definitions and the support for dynamic
refinements in process models.
In this paper, the research question of interest is “What are the benefits of combining goal
modelling with process modelling?”. We are answering this question by focusing on two decades
of experience in using URN in multiple academic
and industrial contexts – based on coarse-grained
statistics from a literature review and the authors’
personal experiences with URN. Furthermore, we
discuss existing and new application areas for
URN to give insight into these benefits. Although
UCM and GRL (or their own ancestors) were used
individually in the 1990’s, their synergical combination only started to appear about two decades
ago. URN enables a rich set of language features,
analysis capabilities, and application areas that
justify its selection for answering this research
question.
This paper contributes a contemporary
overview of the URN language, of its general usage, and of the benefits brought by its
combination of goal modelling with process
modelling along many lines of research. This
knowledge is beneficial to researchers interested
in the synergy between goal and process modelling in general (and not just with URN), and to
practitioners looking for immediate or potential
solutions related to the application domains
discussed in this paper.
This paper is structured as follows. Sect. 2
first provides a historical perspective on the standardization of URN, together with an introduction
to URN’s syntax, semantics, and analysis capabilities, through an illustrative example. Sect. 3
situates URN by providing a trend analysis of various areas where this language has been applied
over the last 25 years. Then, Sections 4 to 9 discuss
in more detail the benefits of using integrated goal
and process models in the domains of goal/process
alignment, regulatory compliance and intelligence,
process adaptation and improvement, value cocreation and service systems, and goal-oriented
process mining, as well as other advanced URNbased modelling techniques, respectively. These
sections present a significant body of literature
reporting on empirical studies in varied application areas, reflecting the experiences of dozens
of researchers and practitioners in applying and
evaluating URN. The benefits of using URN in
these domains, and integrated goal and process
models more generally, are highlighted for each
domain. Sect. 10 presents research challenges and
opportunities for the next decade of URN-based
modelling. The conclusion summarizes the paper.
International Journal of Conceptual Modeling
Vol. 17, No. 2 (2022). DOI:10.18417/emisa.17.2
4 D. Amyot, O. Akhigbe, M. Baslyman, S. Ghanavati, M. Ghasemi, J. Hassine, L. Lessard, et al.
Special Issue on 100 Years of Graphical Business Process Modelling by R. Laue, H. Mayr & B. Thalheim
2 URN Overview
After giving a historical perspective, this section
highlights, with the help of an illustrative example,
the graphical and textual syntaxes of the process
and goal modelling parts of URN. Basic analysis
capabilities based on the URN semantics are also
presented.
2.1 URN History
The history of URN’s building blocks goes back
nearly three decades. The Use Case Maps notation
was first proposed by Buhr and Casselman (1995)
from Carleton University, with a revised version in
a seminal paper by Buhr (1998). However, UCM
was itself based on earlier work by Vigder (1992)
on design slices, under Buhr’s supervision. Slices
helped visualize the behaviour of cross-cutting
concurrent software components. They were extended to support the description of scenarios in
object-oriented and real-time systems and became
timethreads, which were themselves renamed Use
Case Maps around 1995. UCM was then presented
as a notation for “describing, in a high-level way,
how the organizational structure of a complex
system and the emergent behaviour of the system
are intertwined” (Buhr 1998). The notation was
not meant to be executable or analyzable through
automated means.
Around the same period, several researchers
at the University of Toronto developed two modelling notations that would form the basis of GRL,
namely the Non-Functional Requirements (NFR)
Framework (Chung et al. 2000; Mylopoulos et al.
1992), and the i* framework (Yu 1997; E. S. Yu et
al. 2011). The first version of GRL was produced
by Liu and Yu around 20001 and proposed a graphical syntax based mainly on a subset of i* that
focused on intentional elements, links, and actors
(for social modelling), supplemented with additional concepts supporting the kind of propagationbased analysis found in the NFR framework. A
textual syntax and an interchange XML schema
were also proposed.
1 See https://www.cs.toronto.edu/km/GRL/
In 1999, a large telecommunication company
(Nortel), heavily involved in international standardization activities, was looking for a new language
that would enable engineering standardization bodies to describe new wireless telephony services
in an abstract way. Their main problem was that
the sequence diagrams and Message Sequence
Charts used as notations at the time required a
commitment to networking infrastructure that varied substantially across companies and countries.
They became interested in UCMs as they decoupled behaviour from the underlying structure of
components. UCMs were also used in other companies including Mitel, which was also in favor of
standardization. Mitel was using goal modelling
for documenting architectural decisions, and they
suggested combining goal and scenario modelling
as part of the proposal for a new standard, especially to support reflective systems (taking social
aspects into consideration), and for feature personalization. Motorola also supported this way
forward.
Requirements for the new language were first
approved by ITU-T in 2003 as Recommendation
Z.150, which was updated eight years later (ITUT 2011). The URN language description itself
was first standardized in 2008 as Recommendation Z.151 (ITU-T 2018), under the leadership of
Amyot and Mussbacher from the University of
Ottawa. The 2008 URN standard contains a definition of the abstract syntax with an integrated metamodel for goal and process modelling, a graphical
syntax for UCM and GRL, an XML-based interchange format, and modelling elements supporting
the definition of analysis contexts. URN was revised in 2012 mainly to support indicators in GRL,
and to improve the execution mechanism of UCM
models (called traversal). The 2018 version of the
standard now includes a textual syntax for URN,
similar to a programming language.
In the URN standard, the semantics of URN
are specified with a set of detailed requirements
for the execution of UCM models and a detailed
description of GRL analysis algorithms. Furthermore, the execution semantics of UCM models
have been formalized outside the standard with
Enterprise Modelling and Information Systems Architectures
Vol. 17, No. 2 (2022). DOI:10.18417/emisa.17.2
Combining Goal modelling with Business Process modelling 5
Special Issue on 100 Years of Graphical Business Process Modelling by R. Laue, H. Mayr & B. Thalheim
the Lotos process algebra (Amyot and Logrippo
2000) as well as with Abstract State Machines
(Hassine et al. 2005a,b). GRL, on the other hand,
has been formalized by mapping GRL models to
constraint satisfaction problems (Fan et al. 2018).
Other historical details are provided in a previous
survey (Amyot and Mussbacher 2011).
Interestingly, the URN standard did not at first
present UCM as a process modelling language,
but rather as a scenario modelling language. This
changed later once URN started being used heavily
for business process management.
2.2 Graphical Syntax
URN’s graphical syntax is introduced here through
an illustrative example. An Incident Management
System (IMS) is typically used in a hospital to
report and categorize incidents, which may be
clinical (e. g., incorrect medication dosage administered) or not (e. g., a visitor slides and falls on the
floor). An incident is first reported by an observer
(e. g., an employee or a visitor), and then reviewed
by a committee to categorize it. This is important
in order to mitigate such incidents in the future,
both to improve patient safety and avoid potential
legal actions. The review committee must include
a physician in case the incident is of a clinical
nature. As the committee needs to be efficient,
the IMS can provide advanced features to automatically request and find available physicians.
Another feature that can improve effectiveness
is the early detection of suspected duplicate incidents, which can be quickly confirmed by the
original observer.
2.2.1 UCM
Fig. 1 shows a UCM diagram (called a map)
describing this process. Note that the path highlighted in red in the figure represents one specific
scenario supported by the process model; this will
be explained further in Sect. 2.4.
A UCM map allows for the process elements
(start/end points, responsibilities, forks/joins,
stubs, timers, and others presented later in this section) to be allocated to components, which can be
of different natures. A component may be of type
team (e. g., IMS) and represent the system itself
or a part of the system at any level of granularity.
Figure 1: Incident reporting and categorization process in an IMS: Root map in UCM.
International Journal of Conceptual Modeling
Vol. 17, No. 2 (2022). DOI:10.18417/emisa.17.2
6 D. Amyot, O. Akhigbe, M. Baslyman, S. Ghanavati, M. Ghasemi, J. Hassine, L. Lessard, et al.
Special Issue on 100 Years of Graphical Business Process Modelling by R. Laue, H. Mayr & B. Thalheim
(a) Normal review UCM, for dynamic stub Review in Fig. 1.
(b) Clinical review UCM, for dynamic stub Review in Fig. 1.
(c) Physician finding UCM, for stub FindPhysician in Fig. 2b.
Figure 2: Sub-processes, also modelled as UCM maps, bound to their containing stubs.
A component of type actor (e. g., Observer and
ReviewCommittee) describe anyone or anything
that interacts with the system. Components can
also contain other components, and a system component may be deemed a passive component of
type object (e. g., the IMS contains a passive
component representing the incident Database).
A map has at least one start point expressing a precondition and/or triggering event (e. g.,
IncidentObserved), and at least one end point
expressing a postcondition and/or resulting event
(e. g., Thanked and Confirmed). Responsibilities
(e. g., Report Incident) describe the activities
performed along a path. Alternatives are captured
with OR-forks and their two or more forked
paths can be guarded with conditions (e. g., [!DubSusp], where ! represents negation and DubSusp
is a Boolean variable that is true iff a duplicate
incident is suspected). Paths can also merge with
an OR-join , but well-nestedness is not required
by the language. Similarly, concurrency can be
introduced along AND-forks and paths can be
synchronized via AND-joins .
As in many other modelling languages, complex
Enterprise Modelling and Information Systems Architectures
Vol. 17, No. 2 (2022). DOI:10.18417/emisa.17.2
Combining Goal modelling with Business Process modelling 7
Special Issue on 100 Years of Graphical Business Process Modelling by R. Laue, H. Mayr & B. Thalheim
processes can be decomposed into reusable subprocesses. A UCM stub is a container for one
sub-process, often called plug-in map in the UCM
literature. Whereas a normal stub contains only
one plug-in, a dynamic stub contains two or
more plug-ins whose selection is determined at
run-time according to a selection policy composed
of a guarding condition for each contained plug-in.
For example, the Review dynamic stub in Fig. 1
contains two plug-in UCMs, one for the normal
review process in Fig. 2a (guarded by the condition
!Clinical) and one for the clinical review process in
Fig. 2b (guarded by Clinical).
There could be many levels of decomposition.
For example, the stub FindPhysician in Fig. 2b
contains a third level of process decomposition
with a plug-in map (Fig. 2c) that describes how the
IMS finds a physician to participate in the review
committee. This particular process uses a UCM
timer that is reset if the physician accepts the
invitation in a timely way, and that times-out and
loops back otherwise.
There are two important characteristics of UCM
modelling illustrated here. First, new components
and containment relationships can be introduced in
sub-processes. Second, some start points and end
points act as connectors to their parent stub, to ensure continuity of processes across decomposition
levels. The binding between the map in Fig. 2c
and its parent stub (FindPhysician) indicates which
input segment of the stub is connected to which
start point (IN1 with PhysicianNeeded here) and
which output of the stub is connected to which end
point (OUT1 with Found here). Unbound start/end
points describe local triggering/resulting events.
2.2.2 GRL
From a goal modelling perspective, Fig. 3 presents
the GRL view of the IMS and its stakeholders, all
represented with actors . Actors have intentional elements such as goals , softgoals ,
and tasks . A softgoal differs from a goal in
that the former can only be satisfied sufficiently
(i. e., actors do not agree on how to measure the
satisfaction of a softgoal or a measure does not
exist for it at all), while the latter can be satisfied
objectively (i. e., there is a standard, agreed-upon
way to measure the satisfaction of a goal). Be
Efficient is an example of a softgoal, as it is always
possible to be more efficient. Be Legally Compliant is an example of a goal, as it is possible to
determine whether one is legally compliant or not.
Tasks (e. g., Duplicate Detection) often describe
operational activities of solutions but may also
describe characteristics of solutions such as their
speed or resource usage. Intentional elements
can have a relative importance to their containing actor, shown as a value between parentheses.
Similarly, actors can have a relative importance
level as well. A qualitative scale (high, medium,
low, none) or a quantitative scale ([0..100]) can
be used here. Note that large models can span
multiple diagrams.
Intentional elements can be connected to each
other using three types of links:
• Decomposition links of types AND, OR,
or XOR. For example, Duplicate Detection is
decomposed into two mutually exclusive alternatives. The satisfaction of the decomposed
element is the minimum satisfaction of the
decomposing elements for AND, and the maximum for OR and XOR. In XOR’s case, only
one alternative can be satisfied at a time.
• Contribution links with a qualitative or
quantitative weight (in [−100..100]). Fully Automated contributes negatively (−25) to Be Effective and positively (25) to Be Efficient. The
satisfaction of an element is the weighted sum of
products (satisfaction × contribution weight) of
the contributors, divided by 100 and truncated
to [0..100].
• Dependency links . For instance, the
satisfaction of Be Legally Compliant cannot be
higher than that of Support Reviews, upon which
it depends.
Unlike most other goal-oriented languages,
GRL supportsindicators , also called Key Performance Indicators (KPIs), which aim to convert
a measured value observed from the external world
to a unit-less satisfaction level in the [0..100] scale
International Journal of Conceptual Modeling
Vol. 17, No. 2 (2022). DOI:10.18417/emisa.17.2
8 D. Amyot, O. Akhigbe, M. Baslyman, S. Ghanavati, M. Ghasemi, J. Hassine, L. Lessard, et al.
Special Issue on 100 Years of Graphical Business Process Modelling by R. Laue, H. Mayr & B. Thalheim
Figure 3: IMS: Goal model in GRL.
that can be propagated to the rest of the goal model
elements. This conversion uses three indicator
parameters, namely the target (corresponding to a
satisfaction of 100 if the observed value meets it or
does better), the worst case (satisfaction of 0 if met
or worse), and the threshold (satisfaction of 50).
Linear interpolation is used for values in between.
For example, the Cost of Legal Actions per Year
indicator in Fig. 3 could have as parameters (in
Euros): target = 0; threshold = 100,000; and worst
= 1,000,000. In this context, an observed yearly
cost of 550,000 Euros would lead to a satisfaction
of 25, whereas a yearly cost of 50,000 Euros would
result in a satisfaction of 75.
In addition to the quantitative scales discussed
earlier for contribution weights and satisfaction
levels, the URN standard also supports corresponding qualitative scales, shown in Fig. 4. These
labels represent combinations of sufficient/insufficient positive/negative qualitative values. Such
labels are useful in a context where little information is known about quantities, which is often
the case in early requirements engineering activities. Numerical values in quantitative scales are
often inferred from historical data or from experts’
opinions through consensus-building mechanisms
(Akhigbe et al. 2014; Liaskos et al. 2012; Vinay
et al. 2014). There are also methods available
Enterprise Modelling and Information Systems Architectures
Vol. 17, No. 2 (2022). DOI:10.18417/emisa.17.2
Combining Goal modelling with Business Process modelling 9
Special Issue on 100 Years of Graphical Business Process Modelling by R. Laue, H. Mayr & B. Thalheim
to validate GRL models, for example, based on
surveys (Hassine and Amyot 2016), and resolve
conflicts regarding modelled quantities (Hassine
and Amyot 2017). The rest of this paper will use
quantitative values in its examples.
Figure 4: GRL qualitative labels for contribution
weights (top) and satisfaction levels (bottom).
2.2.3 URN Links and Metadata
URN supports user-defined typed links that can be
used to connect any pair of modelling elements.
These are particularly useful to trace elements
from the GRL view to those of the UCM view
in support of rationale documentation. GRL and
UCM elements are annotated with a triangle that
indicates the existence of incoming links or
outgoing links . For example, there is a URN
link of type traces from the GRL task Duplicate
Detection in Fig. 3 to the UCM responsibility
Check Duplicate in Fig. 1.
URN also supports the tagging of model elements with metadata, which are user-defined
name-value pairs. The use of metadata and URN
links enables profiling URN to specific domains,
such as regulations (Ghanavati et al. 2014a). Such
concepts can also be exploited during domainspecific analysis.
2.2.4 Cognitive Effectiveness
In terms of conceptual coverage, effectiveness, and
other such metrics, the UCM and GRL graphical
syntaxes fare well compared to similar languages.
For example, Abrahão et al. (2019) report empirical evidence that the quality of goal models
obtained with GRL is higher than that of i*, and
that GRL is perceived as easier to use and more
useful than i* in several experiments. Mussbacher
and Amyot (2008) also assessed the benefits of
UCM over BPMN and UML Activity Diagrams
in terms of their conceptual coverage of common
workflow patterns.
From a cognitive effectiveness however, both
GRL and UCM suffer from symbol deficits(Genon
et al. 2011; Moody et al. 2010). For example, the
URN standard does not provide a graphical syntax
for i) binding relationships between stubs and
plug-in maps, ii) the details of Boolean conditions,
iii) scenario definitions, as well as iv) UCM’s
many performance-related attributes (workloads
on start points, probabilities on OR-forks, resource
utilisation by responsibilities, etc.) not covered in
this paper. See the work of Petriu et al. (2003) for
more information on performance aspects in UCM
models. On the GRL side, there are also deficits
with the absence of concrete syntax for i) strategies,
ii) indicator definitions, and iii) advanced concepts
such as contribution overrides. URN metadata
is not visualized, and typed URN links are not
visualized beyond their presence, indicated by
triangles (Amyot et al. 2012).
2.3 Textual Syntax
A textual syntax for URN models was recently
standardized in 2018 (ITU-T 2018), and Kumar
and Mussbacher (2018) illustrate its use. This
syntax addresses the symbol deficit issues raised
in the last section, and enables efficient coding
and management of models using standard development environments.
Listing 1 presents an example of the use of the
textual syntax with an excerpt of the GRL model
from Fig. 3, including the two softgoals of the
Committee actor and two tasks of the IMS actor.
The textual syntax allows a short name and a full
name of an element to be defined (e. g., C and
Committee). The short name is used to refer to the
element (e. g., in the two contributions of the first
task of IMS).
Listing 2 presents an example of the use of
the textual syntax for the Clinical review map from
Fig. 2b. At the top, the path is specified, including
the start point ClinicalReview> and end point ReviewDone., responsibilities (e. g., "Get Physician"),
the stub FindPhysician with its plug-in binding
International Journal of Conceptual Modeling
Vol. 17, No. 2 (2022). DOI:10.18417/emisa.17.2
10 D. Amyot, O. Akhigbe, M. Baslyman, S. Ghanavati, M. Ghasemi, J. Hassine, L. Lessard, et al.
Special Issue on 100 Years of Graphical Business Process Modelling by R. Laue, H. Mayr & B. Thalheim
specified in parentheses, as well as the AND-fork
{| with its two branches and the AND-join |}.
1 actor C #" Committee " {
2 softgoal effv #" Be Effective " {
3 importance 50
4 }
5 softgoal efft #" Be Efficient " {
6 importance 50
7 }
8 }
9
10 actor IMS #" IMS " {
11 task fAut #" Fully Automated " {
12 contributesTo C . effv with -25
13 contributesTo C . efft with 25
14 xor decomposes dDet
15 }
16 task dDet #" Duplicate Detection " {
17 ...
18 }
19 ...
20 }
21 ...
Listing 1: IMS: Textual GRL syntax for the IMS goal
model (excerpt).
Connections between path nodes are indicated
with ->. At the bottom, the containment of path
elements and/or components in a component (e. g.,
Review Committee) is specified.
1 map " Clinical review " {
2 ClinicalReview > ->
3 " Get Physicican " ->
4 FindPhysician (
5 " Physician finding ":
6 found = out1 ,
7 PhysicianNeeded = in1 ) ->
8 {|
9 " Analyze Report ".
10 " Analyze Clinical Impact ".
11 |} ->
12 " Categorize Incident " ->
13 ReviewDone .
14
15 actor ReviewCommittee :
16 ClinicalReview ,
17 " Get Physician " ,
18 " Analyze Report " ,
19 " Categorize Incident " ,
20 ReviewDone ,
21 Physician
22 actor Physician :
23 " Analyze Critical Impact "
24 team IMS :
25 FindPhysician
26 }
Listing 2: IMS: Textual UCM syntax for the IMS
Clinical review process.
2.4 URN Modelling and Analysis
On one hand, the URN standard itself does not impose a methodology for creating models. The first
methodology for iteratively creating URN models, including their GRL and UCM views, was
proposed by Liu and Yu (2004), with a focus on
information systems. One good observation here
is that goal elicitation is often easier to do with
managers working at a strategic level whereas
process elicitation is often more effective with
stakeholders closer to the operational level. On
the other hand, the URN standard formalizes modelling concepts and guidelines for analysing GRL
models and UCM models.
The analysis of GRL models is done using
strategies, which are sets of initial satisfaction
values for intentional elements and of observed
values for KPIs, and propagation algorithms that
compute the satisfaction of the other linked elements, their actors, and the entire model. These
algorithms can use as input qualitative satisfactions, quantitative satisfactions, or a mix of both
(Amyot et al. 2010). The jUCMNav tool, a free
Eclipse-based environment for URN modelling
and analysis2 , supports seven such algorithms
(Amyot et al. 2012).
For example, Fig. 5 shows the jUCMNav tool
with the evaluation result of the GRL model in
Fig. 3 for a given strategy, where the With Confirmation and Automatic Request tasks have been
selected, and where the three KPIs have been fed
observed values. The KPIs first compute their
corresponding satisfaction levels as explained earlier, and then the quantitative algorithm used here
propagates these initial satisfaction levels to the
2 https://github.com/JUCMNAV/projetseg-update/wiki
Enterprise Modelling and Information Systems Architectures
Vol. 17, No. 2 (2022). DOI:10.18417/emisa.17.2
Combining Goal modelling with Business Process modelling 11
Special Issue on 100 Years of Graphical Business Process Modelling by R. Laue, H. Mayr & B. Thalheim
Figure 5: GRL model, evaluated for one strategy, in the jUCMNav tool.
other model elements, in a bottom-up way. Colour
feedback is also provided: the greener the better,
and the redder the worse.
Different strategies can be defined and compared, hence supporting "what-if" analysis as well
as trade-off assessments (i. e., which stakeholders
will be satisfied or not).
Analysis of goal models can also be done in
a top-down way, i. e., to produce a GRL strategy
that will optimize some goals given other constraints. For example, Fan et al. (2018) recently
provided arithmetic semantics for GRL models
and automated the generation of mathematical
functions from GRL models, in jUCMNav. These
functions can be fed to optimizers such as CPLEX
(IBM 2019) in order to find an optimal trade-off,
represented in the form of a strategy, given some
optimization objectives and constraints (Anda and
Amyot 2020). On the UCM side, the URN standard provides the concept of scenario definition,
which defines the initial context used to traverse
a UCM model. A scenario definition specifies
which start point(s) are triggered, initial values for
the model’s variables, expected end points reached
at the end, as well as optional pre/post conditions.
A UCM model can include Boolean, Integer,
and enumeration variables used to specify guarding conditions (using a syntax based on Java).
For example, the IMS model used here (Figs. 1
and 2) contains Boolean variables (DupSusp, DupConf, and Clinical) and Integer variables (Attempt
and NumAttempt). UCM responsibilities can also
contain code that updates these variables during
the traversal. For example, Attempt is set to 0 by
responsibility Get Physician and is incremented by
1 by responsibility Ask To Join.
International Journal of Conceptual Modeling
Vol. 17, No. 2 (2022). DOI:10.18417/emisa.17.2
12 D. Amyot, O. Akhigbe, M. Baslyman, S. Ghanavati, M. Ghasemi, J. Hassine, L. Lessard, et al.
Special Issue on 100 Years of Graphical Business Process Modelling by R. Laue, H. Mayr & B. Thalheim
Figure 6: IMS: Traversal result for a UCM scenario
definition, shown as a sequence diagram.
The path highlighted in red in Figs. 1 and 2
results from the traversal of the model for a
scenario definition by jUCMNav. In this definition, the start point is IncidentObserved and
the initial values specify that there is a duplicate
incident suspected (DupSusp=true) but not confirmed (DupConf=false) for a non-clinical incident
(Clinical=false). The traversal detects that the normal review sub-process must be used here (i. e.,
Fig. 2a) and not the clinical review.
The UCM model traversal based on scenario
definitions enables testing the process model and
avoiding regression during modifications. jUCMNav also enables exporting and visualizing traversal results as sequence diagrams. For example,
Fig. 6 corresponds to the flat representation of the
path highlighted in red in our example. Such diagrams are helpful to communicate with stakeholders and systems engineers, and even for deriving
test cases.
3 URN Application Areas
The literature review from (Amyot and Mussbacher 2011) identifies multiple application areas
for URN models, from telecommunication services and business models to healthcare systems.
Using queries on Google Scholar and Scopus,
this section provides an updated perspective with
coarse-grained statistics for important clusters of
application areas.
In order to get an order of magnitude of the
number of URN-related papers published in the
past few decades, we ran the following query on
Google Scholar (including patents but excluding
citations) on March 8, 2021:
"User Requirements Notation" OR "Use Case
Map" OR "Use Case Maps" OR "Goal-oriented Requirement Language" OR "Rec. Z.151" OR "Recommendation Z.151" OR jucmnav.
The acronyms (URN, UCM, and GRL) were
excluded as they led to too many false positives.
However, jUCMNav was included as the main
URN modelling environment. This resulted in
3,730 results, including 70 patents. The first
papers were from 1995, which is the year the "Use
Case Map" expression was coined.
As Google Scholar is known to index many lowquality documents and some grey literature, and as
it indexes the full text of the papers (meaning that
URN could just be mentioned in passing) a similar
query was ran on a multidisciplinary, reputable,
and curated search engine. The results presented
in the remainder of this paper are based on this
Enterprise Modelling and Information Systems Architectures
Vol. 17, No. 2 (2022). DOI:10.18417/emisa.17.2
Combining Goal modelling with Business Process modelling 13
Special Issue on 100 Years of Graphical Business Process Modelling by R. Laue, H. Mayr & B. Thalheim
curated search engine and not the Google Scholar
search. Elsevier’s Scopus3 was selected here
as it offers a flexible query language, it enables
searches on various publication metadata (leading
to higher precision in the selection of papers that
actually focus on URN), and it has good analytics
capabilities. The following query was used:
ALL( "User Requirements Notation" OR "Use
Case Map*" OR "Goal-oriented Requirement Language" OR "Rec. Z.151" OR "Recommendation
Z.151" OR jucmnav )
AND ( EXCLUDE (DOCTYPE , "cr") )
The last part of this query excludes conference
reviews (e. g., prefaces of proceedings). This
query on all metadata (but not on full texts) resulted in 1,466 URN-related papers since 1995,
from 76 countries. A more restricted version of
this query scoped to titles, abstracts, and keywords
(i. e., using TITLE-ABS-KEY instead of ALL) resulted in 305 papers, from 37 countries. These
are papers with core contributions to URN or that
use URN substantially.
Fig. 7 displays the distribution of these papers,
per year. The dotted line (all metadata) uses the
vertical axis on the right of the figure whereas
all six other lines use the one on the left. The
Title/Abstract/Keywords distribution is the one
discussed earlier. The five other lines each represent a category of papers where the query based on
all metadata was restricted (using AND) with additional terms for popular application areas searched
in titles, abstracts, and keywords only.4
• Telecom/Networks (89 papers): Telecommunication features, telephony services, and other
networking protocols and applications collectively represent the original area for which URN
was developed. URN quickly became popular
here, even before its standardization, especially
for validation and verification purposes. However, this area slowly faded away as other areas
became more popular in the last decade.
3 https://www.scopus.com/
4 The Scopus queries and raw numbers are available as a
companion file in this submission.
• Software/Enterprise Architecture (139 papers):
One important area concerns the application of
URN to software architectures and enterprise
architectures, including the reverse-engineering
of existing architectures, trade-off analysis, and
performance analysis. Interest in this area
peaked around 2015 and also seems to slowly
fade away.
• Laws/Regulations (133 papers): Another popular area, more surprising this time, relates
to legal and compliance aspects. Many papers used URN to model laws, regulations,
and policies, often using GRL, together with
mandatory processes in UCM. Analysis often
involved traceability (for change management),
compliance analysis, and performance analytics. Interest in using URN for legal aspects
has increased substantially since 2009, and has
shown the highest level of activity for the past
5 years among the five areas we identified.
• Business/Process/Value (110 papers): For the
past 15 years, and at a sustained frequency,
URN had much success in its application to
business process modelling and management
(often exploiting the UCM view), but also in
business and value modelling (often exploiting
the GRL view).
• Health/Medicine (71 papers): In the last decade,
one increasingly important type of business
area targeted by URN publications relates to
healthcare and medicine, often from a clinical
perspective (e. g., medical procedures), but also
from an administrative perspective, for regulatory compliance, process improvement, or
change management.
Of course, the above assessment is subject to
many limitations as only one search engine was
used, and there likely are relevant papers not included. These categories are also not mutually
exclusive, i. e., one paper could be part of many
categories. Still, the historical trends and proportions are telling and likely representative of reality
in the scientific literature.
International Journal of Conceptual Modeling
Vol. 17, No. 2 (2022). DOI:10.18417/emisa.17.2
14 D. Amyot, O. Akhigbe, M. Baslyman, S. Ghanavati, M. Ghasemi, J. Hassine, L. Lessard, et al.
Special Issue on 100 Years of Graphical Business Process Modelling by R. Laue, H. Mayr & B. Thalheim
Figure 7: Number of URN papers per year, for several categories, based on Scopus queries ran on March 8, 2021.
A closer look at these papers, this time using a
Scopus query that looks for titles, abstracts, and
keywords that cover both UCM and GRL, suggests
that only about 10% of the above papers actually
address a combination of goal and process modelling. This still concerns a significant number of
papers, and many contributions combining goal
modelling in GRL with process modelling in UCM
will be covered in the next sections.
Interestingly, many patents (70 according to
Google Scholar) were also influenced by URN.
Some simply use UCM models to describe scenarios or processes as part of the patent documentation, e. g., see Liscano et al. (2008) and Mears
et al. (2016). Other patents however build their
solutions on top of URN, for example by UCM
modelling as part of a larger method (Dotan et al.
2013), UCM models as input (Weigert 2017), or
GRL models for system configurations (Franchitti
2019). A commercial version of the jUCMNav
tool with additional verification and test generation
capabilities for UCM models was also produced
by Updraft (formerly Uniquesoft, see Weigert et al.
(2019)). This shows that URN is used both in
academia and in industry.
While this section grouped the papers into application areas such as Telecom/Networks and
Health/Medicine, the following six sections offer
an orthogonal grouping into key research areas
focusing on combined goal/process modelling that
is based on the literature review as well as the
authors’ personal experiences.
4 Goal/Process Alignment
Organizational alignment is an area that certainly
benefits from a combination of process and goal
modelling. Whether the existence and definitions
of processes meet organization and stakeholder
goals, and whether goals are properly covered by
such processes are important questions to answer,
especially as both goals and processes evolve.
In addition to the work mentioned in the introduction (Cardoso et al. 2011; Koliadis and
Ghose 2006), other work in that area includes the
approach of Guizzardi and Reis (2015), which
provides an informal and manual method for tracing BPMN processes to Tropos goals in order to
assess alignment. Sousa and Leite (2014) studied
the merging of i* goals and BPMN processes, and
the addition of indicators to trace and measure
Enterprise Modelling and Information Systems Architectures
Vol. 17, No. 2 (2022). DOI:10.18417/emisa.17.2
Combining Goal modelling with Business Process modelling 15
Special Issue on 100 Years of Graphical Business Process Modelling by R. Laue, H. Mayr & B. Thalheim
alignment. More recently, Insfrán et al. (2017)
explored a combination of GRL with BPMN in an
approach that promotes traceability and process
alignment by construction, with several mapping
and prioritization rules. This paper is particularly interesting in its combination of GRL with a
process language other than UCM. Behnam and
Amyot (2013) also explored the selection and configuration of reusable process patterns in UCM
based on GRL goal models evaluated based on a
given organizational context.
Most of these approaches attempt to establish
traceability between goals and processes at construction time. While there are benefits in ensuring
that processes satisfy their objectives when they
are built or selected, such approaches are limited
in their assessment of alignment when the involved
goals and processes evolve over time.
Akhigbe et al. (2016) provided an initial set
of consistency and completeness rules that help
establish and maintain alignment between GRL
goals and UCM processes. Given the importance
of a basic infrastructure to assess alignment, this
work was extended by i) considering new rules
that exploit decomposition structures in these models, ii) automating the verification of these rules
in jUCMNav, iii) providing usability features in
jUCMNav to simplify the tagging of elements,
and iv) providing model transformations for reestablishing alignment in case of rule violations.
Alignment rules are not trivial when balancing
correctness with usability, especially for large
models. Simple rules requiring that all goals be
covered by a process or an activity are simple to
implement, but:
1. Not all rules apply the same way to all goal/process models; they must be selectable by modellers.
2. Not all intentional elements in a goal model are
operationalized as process model elements (indicators being a good example here). Similarly,
not all process elements need a justification
captured in a goal model. Such GRL and UCM
model elements must not be subject to the alignment rules, to avoid false positives.
3. There are just too many elements to align in real
models; some rules must allow taking advantage of model structures (scoped to diagrams,
containers, etc.) to imply coverage.
4. Checking the structural presence or absence
of certain links is simple, but ensuring that the
two aligned model elements are the right one is
difficult.
A set of 18 goal-process alignment rules is
provided for URN5 , importable in jUCMNav, to
address the first three issues discussed above. The
fourth issue remains to be addressed at this time.
The rules are formalized as OCL constraints that
can be checked selectively against URN models,
using the jUCMNav’s constraint selection and
verification feature (Amyot and Yan 2008).
These rules exploit a specific metadata
(Traces=No) that indicates, when used to annotate a model element, that this element is
outside the scope of the rules. jUCMNav now includes a button to add and remove this annotation
in one click. The rules check the presence of a
particular type of URN link (see Sect. 2.2.3) of
type traces. Often, these rules come in pairs, for
instance:
• Each GRL actor must have a Traces link
to a UCM component, unless tagged with
Traces=No.
• Traces links from a GRL actor must only be to
a UCM component.
The first one looks for the existence of a UCM
component aligned with that GRL actor, whereas
the second one restricts the type of links from
GRL actors. A similar pair of rules exist from
UCM components to GRL actors.
Some rules are also meant to be mutually exclusive, for instance:
• Each GRL intentional element must have a
Traces link to a UCM map or responsibility,
unless tagged with Traces=No.
5 Available at http://bit.ly/URN-alignment-rules
International Journal of Conceptual Modeling
Vol. 17, No. 2 (2022). DOI:10.18417/emisa.17.2
16 D. Amyot, O. Akhigbe, M. Baslyman, S. Ghanavati, M. Ghasemi, J. Hassine, L. Lessard, et al.
Special Issue on 100 Years of Graphical Business Process Modelling by R. Laue, H. Mayr & B. Thalheim
• Each GRL task must have a Traces link to a
UCM map or responsibility, unless tagged with
Traces=No.
The first one requires the alignment of all types of
GRL intentional elements (goals, softgoals, and
tasks) with UCM process elements whereas the
second one restricts this alignment to GRL tasks
only.
To minimize the number of Traces links that
must be manually created by the modeller, some
rules infer alignment of finer-grained model elements from the alignment of more coarse-grained
elements. In this context, alignment can take advantage of containment relationships (intentional
element to actors in GRL, or path elements to
component or maps in UCM) and refinement
relationships (e. g., decomposition in GRL and
stub/plug-in relationships in UCM). Such rules
include:
• Each UCM component (or one of its parents)
must have a Traces link from a GRL actor,
unless tagged with Traces=No.
• Each GRL intentional element (or one of its
decomposed parents) must have a Traces link
to a UCM map or responsibility, unless tagged
with Traces=No.
The above rules are recursive in order to support
multiple levels of containment and refinement.
Such powerful rules must be used with care as
they will cover many sub-elements all at once,
including some that might be unjustified.
Fig. 8 shows the result of applying 14 of the 18
alignment rules on the IMS example from Sect. 2.
Seven instances of violations of four different
rules are reported, together with the rule name
and the location in the model. Right-clicking on a
warning also enables a contextual menu to offer
three typical types of fixes: i) delete the element,
ii) create a new element (of the type expected
by the rule) with the expected traces link, and
iii) ignore the element (by adding the Traces=No
metadata).
In this example, the Hospital GRL actor does not
have a corresponding UCM component. This actor
should likely be ignored as it is not involved in the
process. The GRL task Fully Automated should
likely be ignored as well (as it was not selected
for this process), whereas the GRL task Automatic
Reminder might be kept but with corresponding
(and linked) improvements to the UCM process.
The UCM component Physician is also not aligned,
and one solution would be to augment the goal
model with a new corresponding (and linked)
actor, with its goals. Such decisions need to be
done for the other warnings as well, and alignment
reassessed after these modifications. Using a
different selection of rules would also lead to a
different set of warnings.
There are clear benefits here in not only establishing alignment when one view is generated
from the other, but in continuously maintaining
alignment when any of the goal or process views
is modified. Only the rules that make sense need
to be checked, and only on the modelling elements
that require it. Recursive rules that take advantage of structural aspects such as refinement and
decomposition also help minimize manual labour
for large models. Tool support where both views
are integrated, links can be added at construction
time, and rule violations can be fixed in one user
action also help with the usability of the alignment
maintenance.
5 Regulatory Compliance and Intelligence
Regulatory compliance aims to ensure that the
business operations of an organization are aligned
with relevant laws, regulations, policies, and other
such obligations. These business operations can
be modelled with business processes, but most
regulatory aspects, which are more often based on
outcomes than on prescribed behaviour, cannot all
be modelled with processes. A recent survey by
Hashmi et al. (2018) indicates two important types
of regulatory compliance: at design time (which
often implies the use of models and traceability)
and at run time (which focuses on checking behavioural instances against expected obligations
and processes).
Enterprise Modelling and Information Systems Architectures
Vol. 17, No. 2 (2022). DOI:10.18417/emisa.17.2
Combining Goal modelling with Business Process modelling 17
Special Issue on 100 Years of Graphical Business Process Modelling by R. Laue, H. Mayr & B. Thalheim
Regulatory Compliance Software Engineering,
which focuses on the development of systematic approaches to build and maintain legally-compliant
software systems, is an area of research that benefited from a combination of goal and process
modelling. In the last decade, with the introduction of new regulations such as HIPAA (U.S.
Government 1996) for US healthcare and GDPR
(Intersoft Consulting 2016) for EU privacy, goal
modelling languages, and especially GRL, have
been extended to help model regulations in the
same way as software and business requirements
(Ghanavati 2013; Ghanavati et al. 2014a,b) or to
help evaluate and establish regulatory compliance.
These approaches build on the idea that modelling regulations and legal requirements in the
same notation as other types of requirements will
help improve the understanding of regulations,
manage compliance with regulations in a systematic way, and identify non-compliant instances in
software and processes. Akhigbe et al. (2019)
recently reported on a literature review of methods used for legal and regulatory compliance and
showed that goal modelling methods, such as
those based on GRL, offered more benefits for
modelling compliance tasks in comparison to nongoal-oriented methods, such as logic-based ones.
This is in part because the goal-oriented methods
can be used to model both the intent and the structure of laws and regulations, whereas the other
methods focus mainly on their intent. Although
regulations generally describe only what the intent
of compliance is (leaving some freedom for organizations to decide how best to comply in their
context), on a few occasions, they also prescribe
operational aspects, i. e., high-level processes. To
capture such operational aspects, researchers also
proposed including UCMs in the modelling of regulations and other legal concepts, with traceability
links between the GRL and UCM views.
Ghanavati (2013) and Ghanavati et al. (2014a,b)
developed a domain-specific version of URN (i. e.,
a profile) within the Legal-URN framework, which
models legal requirements with Legal-GRL and
existing legal processes with Legal-UCM. Fig. 9
shows an overview of Legal-URN. Legal-URN
includes three layers where organizational processes/procedures and policies are modelled in UCM
and GRL respectively, and where regulations are
modelled with Legal-GRL and Legal-UCM. Organizational models (for instance, similar to the IMS
example of Sect. 2) and legal models are connected
to each other via various types of traceability links.
In Legal-URN, the top layer (i. e., procedures &
policies and multiple regulations) helps track the
models to their source documents. The second
layer, with organizational GRL and Legal-GRL
Figure 8: Warnings reported by jUCMNav for violated alignment rules, with potential fixes.
International Journal of Conceptual Modeling
Vol. 17, No. 2 (2022). DOI:10.18417/emisa.17.2
18 D. Amyot, O. Akhigbe, M. Baslyman, S. Ghanavati, M. Ghasemi, J. Hassine, L. Lessard, et al.
Special Issue on 100 Years of Graphical Business Process Modelling by R. Laue, H. Mayr & B. Thalheim
models, supports quantitative and qualitative legal
compliance analysis. The last layer, composed
of organizational UCM and Legal-UCM models,
helps ensure business process compliance. LegalURN enables evaluating legal compliance through
GRL strategies using tailored qualitative and quantitative propagation algorithms that exploit the new
profile information.
Legal-GRL introduces new types of intentional
elements for GRL by adding stereotypes (described as URN metadata) to softgoals and goals.
These stereotypes are based on deontic notions of
legal requirements (i. e., permissions and obligations). Softgoals and goals in Legal-GRL are annotated with «Permission» and «Obligation» stereotype metadata. Obligations are intentional elements that are mandatory by law (“shall”, “must”,
and “will” are usually used to indicate obligations
in legal texts), whereas permissions are optional legal requirements (often indicated by the presence
of “may” or “might”). Legal-GRL also includes
other stereotyped intentional elements to illustrate
a «Precondition», an «Exception», or a crossreference «XRef», together with consequences of
non-compliance in legal requirements (specified
using contribution links).
The Legal-URN profile is implemented in
jUCMNav and includes 18 OCL rules used to
verify the well-formedness of Legal-GRL models
and 5 OCL rules to identify non-compliant instances between Legal-GRL and GRL. The LegalURN profile also benefits from the relationship
between GRL and UCM in that the Legal-GRL
intentional elements and their links can be traced
to components, stubs, and processes in the LegalUCM view. To show the traceability between the
two types of models, URN links are used. This
traceability helps propagate the satisfaction values
of intentional elements in Legal-GRL models to
components in Legal-UCM models. Traceability
of the models to laws/regulations and to organizational policies is handled through third-party
requirements management tools, where models
can be imported and modifications to linked models and texts tracked (Rahman and Amyot 2014).
Legal-URN hence benefits process modelling by
better ensuring, at design time, that organization
processes comply with models of laws (involving
both goals and processes) and maintain compliance when changes occur to laws or organization
processes and objectives.
The Variability-Aware Legal-GRL Framework
(Sartoli et al. 2020a,b) is another example of a
domain-specific profile of URN that extends GRL
and Legal-GRL. This framework helps to analyze
compliance of software systems to multiple regulations in the face of run-time variability and
environmental changes and provides adaptive privacy. Adaptive privacy is defined by Sartoli et
al. (2020a,b) as “capabilities to monitor information transmission to allow response to contextual
changes”. In this framework, Legal-GRL’s intentional elements are combined with Privacy-Aware
Role Based Access Control (P-RBAC) elements
(Ni et al. 2010) through a set of stereotypes (described as URN metadata) to ensure capturing new
constructs such as purpose, role, and context, and
support evaluating a legal model based on different
environmental variables. The Variability-Aware
Legal-GRL profile includes three profiled types
of actors (i. e., ≪Compliance Actor≫, ≪Data
Requester≫, and ≪Legal Actor≫). Permission
intentional elements from Legal-GRL are connected to ≪context≫ and ≪ purpose≫ GRL tasks
through decomposition links. ≪Capability≫ elements are represented as tasks and, finally, ≪data
object≫s are represented as resources. This framework benefits process modelling by helping processes better comply dynamically with multiple
regulations.
Legal-UCM has been also used to help model
ambiguities defined by an Ambiguity Taxonomy
(Massey et al. 2014) in legal statements (Massey et
al. 2017). In this work, UCM stubs are used and
stereotyped to cover six categories of ambiguities:
lexical, syntactic, semantic, vagueness, incompleteness, and referential. In this work, the legal
requirements are categorized as non-ambiguous
or ambiguous. If a requirement is non-ambiguous,
it is treated similar to system requirements and
is modelled as-is. Otherwise, if a requirement
is ambiguous, a stub is added in the UCM path
Enterprise Modelling and Information Systems Architectures
Vol. 17, No. 2 (2022). DOI:10.18417/emisa.17.2
Combining Goal modelling with Business Process modelling 19
Special Issue on 100 Years of Graphical Business Process Modelling by R. Laue, H. Mayr & B. Thalheim
Figure 9: Overview of the Legal-URN Profile.
and the type of ambiguity is documented in the
stub. Such augmented process models benefit
legal experts as a means to discuss ambiguities
and resolve them.
Whereas the Legal-URN work mainly focuses
on design-time compliance, a combination of goal
and process modelling was also investigated for a
run-time compliance perspective. The particularity here is the focus on GRL not only to specify
the regulations and policies with which processes
need to comply, but on the use of GRL indicators
to measure this level of compliance in a non-binary
way (Shamsaei et al. 2010). The satisfaction levels
of indicators and other goal model elements are
also exported to a Business Intelligence tool to
support compliance analytics. This led to another
compliance profile for URN and an Indicatorbased Policy Compliance Framework (Shamsaei
2012) with a core focus on aviation and transportation regulations that benefits organizations
in detecting hot spots in processes, captured with
UCM, regarding run-time regulatory compliance.
Further building on this work, Akhigbe (2018)
proposed another URN profile supporting the
Goal-oriented Regulatory Intelligence Method
(GoRIM), this time to help regulators assess
whether regulations themselves achieve their desired objectives and accomplish their intended
outcomes. GoRIM is a method that facilitates
evidence-based decision-making through modelling and data analytics using advanced Business
Intelligence tools, a domain referred to as Regulatory Intelligence (Akhigbe et al. 2017). GRL
enables modelling three complementary views
(regulations, regulatory initiatives used to administer the regulations, and intended social outcomes)
with the same language and enables robust analysis within each view and across views. A data
analytics tool is then used to explore and analyze
data derived from goal models. A combination
of goal and process modelling in methods such
as GoRIM provides benefits to regulators in the
form of an intelligent approach to regulatory management, as well as a feedback loop in the use
of data for continuous monitoring, assessment,
and reporting on efficiency and effectiveness of
regulations.
6 Process Adaptation and Improvement
The URN standard does not specify how the combined analysis of GRL and UCM views can be
done. Early work from Roy et al. (2006), implemented in jUCMNav, enables the automatic
creation of UCM Integer variables for each intentional element in the GRL view. Each variable
International Journal of Conceptual Modeling
Vol. 17, No. 2 (2022). DOI:10.18417/emisa.17.2
20 D. Amyot, O. Akhigbe, M. Baslyman, S. Ghanavati, M. Ghasemi, J. Hassine, L. Lessard, et al.
Special Issue on 100 Years of Graphical Business Process Modelling by R. Laue, H. Mayr & B. Thalheim
represents the satisfaction level of its corresponding GRL element. The benefit of this approach
is that each view can dynamically influence the
other. Satisfaction levels can be used in UCM
conditions (influencing the traversal of the UCM
model) and in turn traversed responsibilities can
include code that modifies the satisfaction level of
goals and other elements. This provides a limited
form of process adaptation that was exploited
for the engineering of Next Generation Networks
services (Amyot et al. 2008).
This synergy was further extended with support
for indicators and a rule-based reasoning engine
(DROOLS) for monitoring context in the work of
Vrbaski et al. (2012) on their CARGO approach.
The benefits here include the modelling of the
feedback loops of self-adaptive socio-technical
systems, reasoning based on a combination of
rules and goal analysis at run-time, and an execution/simulation environment for context-based
process adaptation.
The above process adaptation strategies focus
mainly on run-time adaptation. However, there
is also a need for adapting external processes to
the context of a specific organization, especially
in a technology adoption scenario. For example, a
hospital always strives for improving its services
and managing its business and clinical processes
for optimal allocation of resources. Technologybased solutions aim to enable improvements and
optimization of business processes and delivered
care services. However, introducing new technological solutions (for instance, the Incident Management System of Sect. 2) in a large organization
such as a hospital, is not a straightforward task
as it may cause existing business and clinical
processes to be changed. This may in turn disrupt current practices and increase user resistance
toward those changes. Hence, the impact of the introduced technology and its dependent processes
on stakeholder/organization goals, business objectives, and predefined criteria have to be illustrated
and evaluated systematically.
Kuziemsky et al. (2010) addressed this issue
by using URN to show the possible impact of
healthcare information systems on operational
business processes and business goals, following
five steps:
• Develop a GRL model and define stakeholders
and their dependencies;
• Assign KPIs to each goal in order to measure
its satisfaction level;
• Develop a UCM model to specify existing processes, in order to identify the impact point of
the new system with regards to tasks, goals, and
KPIs;
• Analyze two cases: the current situation with
changes introduced by the new system, and the
current situation without the proposed system
with respect to the tasks and goal;
• Assess the alternatives continuously.
Kuziemsky et al. (2010) applied their framework
on a palliative care process, with observed benefits
regarding the framework’s effectiveness in identifying and prioritizing indicators to decide whether
to implement or resist the proposed changes.
Along the same line, the Activity-based Process Integration (AbPI) framework defines a URN
profile for a goal-driven, analysis-enabled process
integration framework in healthcare (Baslyman
et al. 2017a,b). The aim of this framework is
to integrate activities of new processes, such as
technology-dependent processes, incrementally
into existing business and clinical processes of
an organization, to minimize disruption and keep
existing goals satisfied or see them improved.
The benefit of AbPI’s incremental integration
is to capture and evaluate the impact of each newly
integrated activity on the stakeholder goals so that
the best process integration alternative is selected.
The input to the framework (see Fig. 10) is a
set of Process-Goal Models (PGMs) that each
contains a GRL goal model, many UCM process
models, and a main concern. PGMs are used
to capture current processes and new processes
to introduce, with their respective goal models.
The output of AbPI is a collection of optimal
PGMs (usually with only one element) that satisfy
stakeholder goals, performance objectives, and
predefined criteria. Interactive integration and
Enterprise Modelling and Information Systems Architectures
Vol. 17, No. 2 (2022). DOI:10.18417/emisa.17.2
Combining Goal modelling with Business Process modelling 21
Special Issue on 100 Years of Graphical Business Process Modelling by R. Laue, H. Mayr & B. Thalheim
Figure 10: Activity-based Process Integration framework.
evaluation methods are used in between to explore
and transform the PGMs.
Each activity has a relationship to processes
(current, new, or integrated), relationships to other
activities of current processes (to add to, combine
with, replace, or eliminate other activities), and a
relationship with the role that performs an activity
(add, change, or eliminate the role). In addition,
each activity can contribute positively or negatively to the satisfaction level of goals or change
the satisfaction value of goals. These types of
relationships were implemented using URN typed
links in a URN profile for AbPI. KPIs are used
to monitor the performance of process activities
of processes and define the impact on the goals.
The profile includes OCL constraints to ensure
completeness and consistency of the PGM models.
AbPI was also used in combination with the
Lean change management approach to support
effective integration of technology into healthcare
practices (Baslyman et al. 2019). The philosophy
of Lean is to increase productivity while minimizing cost and eliminating waste. Lean also
uses quantitative measurements to constrain the
design of solutions. AbPI can be used in a way
that supports the Lean approach by providing essential information systematically such as i) the
prioritization of goals based on customer values,
ii) guidance for the integration of a new process
into current processes based on added-value and
non-adding-value activities, and iii) quantitative
measurements with analysis to reason about integration alternatives and select the best one that
meets predefined criteria. An interesting benefit
here is that goal models help formalize social
concerns and measures during the Lean-oriented
change management of processes.
7 Value Co-creation and Service Systems
Service engineering has been core to the development of URN (Amyot et al. 2008). It is also
a domain in which process modelling and goal
modelling have been widely used in the past
two decades, but often separately. One of the
most well-known, process-based modelling notation for designing and managing services is
International Journal of Conceptual Modeling
Vol. 17, No. 2 (2022). DOI:10.18417/emisa.17.2
22 D. Amyot, O. Akhigbe, M. Baslyman, S. Ghanavati, M. Ghasemi, J. Hassine, L. Lessard, et al.
Special Issue on 100 Years of Graphical Business Process Modelling by R. Laue, H. Mayr & B. Thalheim
service blueprinting (Bitner et al. 2008). Service
blueprinting illustrates the sequential processes
composing a service in terms of physical evidence (e. g., parking), customer action (e. g., give
bags), onstage contact activities (e. g., greet), and
backstage activities and support processes, including technological systems. Extensions to service
blueprinting have added graphical modelling techniques to express actors’ goals, among others, in
order to better support the analysis of varied service experiences (Patrício et al. 2018). A number
of propositions have also used goal modelling to
support service analysis and design. For example,
i* has been successfully applied to model recurring business models and patterns in business
services (Lo and Yu 2007). It was also adapted
to create a modelling technique for expressing the
processes by which service actors collaborate to
create valued service outcomes, leading to value
co-creation modelling (Lessard 2015b).
The interest in modelling goals for service design and engineering reflects the importance of
analyzing multiple actors’ perception of the value
generated through the service process (Lessard
2015b). Indeed, while a service expressed as a
process model would end once the service has
been delivered or received, its true endpoint should
be the determination of value by its stakeholders
once they have integrated it (or not) as a resource
within their environment (Lessard et al. 2020).
Moreover, goal modelling notations are able to
express concepts that are key to the engineering of
service systems beyond activities, including agent
roles and dependencies in a service network and
metrics and indicators of value (Lessard 2015a).
URN enables an integrated analysis of a service
system across these dimensions through GRL elements and linked UCM processes. This has been
demonstrated through the creation of a lightweight
GRL profile and accompanying heuristics for
modelling service systems, used to elicit requirements for a telemonitoring service (Lessard et al.
2020). Results showed that requirements identified through the use of the lightweight GRL profile
are more complete than comparable requirements
proposed in the literature, identifying additional
requirements related to the access and use of the
telemonitoring solution. Building on the example of the IMS illustrated in Figs. 1 to 3, such an
application of URN would enable modelling an incident management service in terms of the impact
of reviewed incident reports on the hospital, the
committee, and other stakeholders’ perception of
value. This could uncover other important system
requirements regarding, for example, information
access rights provided to patients about incidents.
In summary, the integration of goal and process
modelling has a number of benefits for service
system engineering:
• The ability to define and measure a service’s
value from varied stakeholder perspectives, in
relation to the processes generating service
outputs and outcomes.
• The ability to redefine the endpoint of a service
process at the moment when value is determined
by stakeholders, in order to capture activities
that impact the perceived success of the service.
• The ability to support the identification of additional system requirements that allow organizational actors to integrate a service’s outputs
and outcomes in a manner that is beneficial to
them.
8 Goal-oriented Process Mining
Process modelling is often done for future to-be
processes. However, there is also a need to model
as-is processes in organizations, ideally based
on evidence rather than on biased perceptions of
what these processes currently are. Process mining is an evidence-based approach to turn event
logs into valuable insights about processes (Van
Der Aalst 2016). In particular, process mining
exploits event logs to discover real processes performed in organizations, enabling modellers to
(re)design and improve process models. Process
mining activities are categorized into three main
types: i) Process Discovery, where a workflow
model is inferred from the event logs; ii) Conformance Checking, where the real processes and
their instances are compared with the desired or
Enterprise Modelling and Information Systems Architectures
Vol. 17, No. 2 (2022). DOI:10.18417/emisa.17.2
Combining Goal modelling with Business Process modelling 23
Special Issue on 100 Years of Graphical Business Process Modelling by R. Laue, H. Mayr & B. Thalheim
prescribed models to find and analyze deviations;
and iii) Process Enhancement, where the desired
data is used to improve or/and extend an existing
process model (Van Der Aalst 2016).
The main input of conventional process mining
activities is the event log, resulting from the execution of processes. A log minimally includes, for
each event, an instance/case identifier, the event
name, and timestamps, optionally with other attributes. Process mining techniques have been
growing into an activity-focused approach that
does not typically consider goals pursued by individual cases and satisfaction levels for different
stakeholders’ goals (Ghasemi and Amyot 2020).
This situation usually threatens the rationality behind the models discovered by process mining
methods and often results in unstructured and
hard-to-read “spaghetti-like” process models. Although such complex models can reflect reality,
they cannot distinguish the traces misaligned with
goals from goal-aligned ones (Pesic and Van der
Aalst 2006). Such problems, especially in flexible
environments that allow many alternatives during
process execution, are to be dealt with by process
mining practitioners.
A literature review of the intersection of goal
modelling and process mining (Ghasemi and
Amyot 2020) showed that although both research
areas are growing, few studies are conducted at
their intersection. Synergetic effects achievable
by combining goal modelling and process mining
can bring important benefits by augmenting the
precision, rationality, and interpretability of the
discovered process models and eventually improve
the satisfaction of process stakeholders.
To address the above research gap, a new approach called Goal-oriented Processes Mining
(GoPM) was recently proposed by Ghasemi and
Amyot (2019b). GoPM, in general, is a process
mining approach concerned with both the sequencing of activities and processes’ goals and satisfaction indicators. In particular, a goal-oriented
process enhancement and discovery (GoPED)
method has been proposed within the GoPM context. GoPED’s log pre-processing step (blue box
on the left of Fig. 11) enhances the input event log
by extracting traces from sequences of interleaved
events and by adding satisfaction levels of corresponding goals and actors to the different traces as
additional attributes (Ghasemi and Amyot 2019a).
GoPED then exploits these new attributes (red box
on the right of Fig. 11) to select the traces that
either meet or do not meet certain goal criteria
(Ghasemi and Amyot 2019b), before feeding the
filtered log to conventional process mining tool
and generate a corresponding process model (in
BPMN, Petri Nets, UCM, etc.).
For example, for the IMS example of Sect. 2, a
trace of activities may reduce legal costs and help
satisfy the goal Legal Actions Avoided but may end
up with an incorrect assessment of incidents that
leads to a higher number of incidents, negatively
impacting the goal High Patient Safety. Other
traces could lead to opposite impacts. GoPED
exploits goal models to manage such conflicting
goals and to support trade-off analysis. Such
analysis helps find, from existing evidence, the
process model promising to satisfy the conflicting goals with a predefined level. In contrast to
conventional process mining that uses the whole
event logs to discover a model, GoPM aims to use
only a subset of the event log that has satisfied
predefined goal-related criteria. With GoPED,
good past experiences will be found within the log
to be used as a basis for discovering good models
for the future and bad experiences will be found
to be avoided.
The goodness of traces (evidence) and models
is defined using three categories of goal-related
criteria: i) satisfaction of individual cases in terms
of some goals (case perspective), ii) overall satisfaction of some goals rather than every single
case (goal perspective), and iii) a comprehensive
satisfaction level for all goals over all cases (organization perspective). Finding the subset of
cases that satisfy a minimum satisfaction level
from each of the above perspectives is not trivial.
This is because adding a trace to the subset might
help satisfy the criterion related to one goal and,
at the same time, might harm satisfying the criterion related to other goals. In GoPED, Ghasemi
and Amyot (2019b) proposed three algorithms
International Journal of Conceptual Modeling
Vol. 17, No. 2 (2022). DOI:10.18417/emisa.17.2
24 D. Amyot, O. Akhigbe, M. Baslyman, S. Ghanavati, M. Ghasemi, J. Hassine, L. Lessard, et al.
Special Issue on 100 Years of Graphical Business Process Modelling by R. Laue, H. Mayr & B. Thalheim
(two of which require constraint-based optimization) for the selection of cases out of enhanced
logs from the three above viewpoints. Recent
experiments demonstrated their feasibility and
scalability. GoPM is expected to guide process
mining approaches towards specific goal-related
properties of interest.
The benefits of goal-oriented process mining
hence extend beyond discovering process models
from evidence. Adding goal models in a process
mining context enables GoPM to discover simpler and understandable process models that meet
measurable goal-oriented criteria. In turn, these
models can help organizations understand and
promote good practices, and help people avoid
process options that lead to underperforming tradeoffs.
9 Other Advanced Modelling Techniques
This section highlights other advanced modelling
and analysis techniques with features, aspects,
slicing, and concepts from social sciences. They
all combine GRL with UCM in a way that complements the previous approaches, and they also
benefit from tool support in jUCMNav.
9.1 URN and Feature Models
In jUCMNav, URN modelling was augmented
with feature modelling to benefit from a richer
description of relationships among features (with
a correspondence to tasks in GRL) and a combined evaluation of feature and goal models with
a single evaluation mechanism (Liu et al. 2014).
For example, the goal model for the IMS actor in
Fig. 3 shows the relationships among tasks with
decompositions in addition to the impacts of those
tasks on high-level goals. However, the modelling
of relationships among tasks is limited in goal
models and distracts from the core purpose of
goal models, i. e., understanding the impact of
tasks. A feature model offers optional, mandatory,
requires, and conflicts relationships in addition
to decomposition with OR and XOR groups. To
benefit from improved separation of concerns,
relationships among features (i. e., tasks) are modelled more expressively with feature models, the
goal model then can focus on the impacts of tasks,
and process (workflow) models describe tasks in
more detail. To unify the evaluation, a feature
model is considered a special case of a goal model.
Optional and mandatory relationships are interpreted as contributions with specific contribution
values that are aligned with the semantics of feature models. Requires and conflicts relationships
are integrated into URN with OCL constraints.
Building on the combined modelling of features and goals and a survey about reusability in
goal modelling demonstrating gaps in the support
of reuse in goal models (Duran and Mussbacher
2019), goals models with relative contribution
values were then applied to decide which features
of a reusable artifact are to be selected to satisfy
high-level goals (Duran 2018). While many kinds
of reusable artifacts exist with clearly defined interfaces, these interfaces focus on structural and
behavioural aspects and not the reasons why an
artifact should be reused. These reasons are captured with goal models in a new kind of interface
for reusable artifacts (Duran 2018). However,
constraints are imposed on tasks in (GRL) goal
Figure 11: Overview of goal-oriented process mining, including the GoPED method.
Enterprise Modelling and Information Systems Architectures
Vol. 17, No. 2 (2022). DOI:10.18417/emisa.17.2
Combining Goal modelling with Business Process modelling 25
Special Issue on 100 Years of Graphical Business Process Modelling by R. Laue, H. Mayr & B. Thalheim
models not only by feature models but also by
(UCM) process models. While a feature model
describes which tasks in a goal model may appear
together in a system, the process model describes
additional sequencing constraints on the tasks
in the goal model that have to be taken into account during goal model evaluation (Duran and
Mussbacher 2016). For example, tasks that always appear on alternative branches in the process
model cannot impact high-level goals at the same
time even though both tasks are selected in the
feature model for the system. In this case, the
accuracy of goal model analysis benefits from
the integration of goal models with feature and
process models. These analysis capabilities were
then further extended with support for top-down
evaluation of goal models (Duran and Mussbacher
2018) that determines the best reusable options
(features/business processes) given desired satisfactions of high-level goals. Issues of computational complexity common in top-down evaluation
is tackled by exploiting the modularization provided by reuse hierarchies.
9.2 Aspect-oriented URN
Aspect-oriented modelling allows crosscutting
concerns to be encapsulated in aspects, which are
then composed to produce an overall model. A
crosscutting concern affects many different elements in a model (e. g., an authentication concern
needs to be applied to each interaction of a user
with the system to ensure security). Aspects generally improve separation of concerns and reusability
of modelling artefacts (Pinciroli et al. 2020).
AoURN (Mussbacher 2010), an aspect-oriented
extension of URN, further enhanced separation
of concerns for stakeholders, use cases/processes,
and non-functional requirements/system qualities.
Integrated goal and process models benefit from
AoURN as aspect-oriented techniques are applied
to both models (Mussbacher et al. 2010a) and
semantic equivalences are exploited in goal and
process models when composing aspectual URN
models (Mussbacher et al. 2013). Aspect-oriented
techniques require a composition specification that
describes where an aspectual model is to be applied. When a model is changed, the composition
specifications often have to be changed, too, so
that the aspects are still applied correctly. The
semantic-based approach of AoURN allows refactoring operations to be performed on an AoURN
model without the need to update composition
specifications, i. e., AoURN is refactoring-safe
(Mussbacher et al. 2009). Furthermore, semanticbased interactions among aspects can be detected
(Mussbacher et al. 2010b). AoURN is also the
basis of an approach based on Software Product
Lines that captures features, goals, and processes
in a unified framework to reason about stakeholders’ needs and perform trade-off analyses while
considering undesirable interactions that are not
obvious from the feature model (Mussbacher et al.
2012).
Fig. 12 illustrates the application of an Authentication aspect to the IMS example. The Security
Server checks whether a User is already authenticated or blocked, and if not asks the User to enter
credentials. If the credentials are correctly entered
within three attempts, authentication is successful. The composition specification identifies the
responsibilities of the Observer as the target of
the Authentication aspect. Although not shown in
the figure, the complete composition specification
applies the aspect also to the Review Committee
and the Physician. When the aspect is applied to
the IMS, the Security softgoal with an impact from
the Authentication task is added to the goal model
of the IMS and an aspect stub is added before
each responsibility of the Observer in the IMS root
map, clearly demonstrating the crosscutting nature
of the aspect.
9.3 URN and Slicing
Similar to aspect-oriented techniques that seek
to handle complexity through separation of concerns, slicing techniques help analysts understand
complex models by identifying relevant model
elements related to a given element of interest
(Androutsopoulos et al. 2013). As UCM process
models evolve and grow over time, they rapidly
International Journal of Conceptual Modeling
Vol. 17, No. 2 (2022). DOI:10.18417/emisa.17.2
26 D. Amyot, O. Akhigbe, M. Baslyman, S. Ghanavati, M. Ghasemi, J. Hassine, L. Lessard, et al.
Special Issue on 100 Years of Graphical Business Process Modelling by R. Laue, H. Mayr & B. Thalheim
(a) Authentication Aspect (UCM).
(b) Composition specification for Authentication Aspect (UCM). (c) Authentication Aspect (Goals).
(d) Authentication Aspect applied to IMS root map.
Figure 12: Example of Aspect-oriented User Requirements Notation (AoURN) applied to IMS.
become hard to understand and to maintain. Indeed, a UCM model may integrate dozens of
maps that are themselves composed of hundreds
of constructs, which make such amount of information humanly unmanageable even by URN
experts. To help analysts understand large and
complex UCM specifications prior to performing a maintenance task, Binalialhag et al. (2018)
proposed a static slicing approach for the UCM
notation, implemented in the jUCMNav tool. The
technique aims to help requirements engineers
reduce UCM specifications according to a slicing
criterion (composed of a UCM construct and a
variable) of interest. The slicing process consists
of a backward traversal of UCM path constructs
and tracking dependencies to identify path elements within a UCM model that may impact the
slicing criterion. The proposed approach produces
both closure and reduced slices. A closure slice is
displayed by marking the relevant constructs and
paths within the original UCM, while a reduced
slice is a new executable UCM model obtained
after the removal of irrelevant constructs, paths,
components, and scenarios.
Fig. 13 illustrates the application of static slicing to the root map of the IMS example, with
Enterprise Modelling and Information Systems Architectures
Vol. 17, No. 2 (2022). DOI:10.18417/emisa.17.2
Combining Goal modelling with Business Process modelling 27
Special Issue on 100 Years of Graphical Business Process Modelling by R. Laue, H. Mayr & B. Thalheim
responsibility Store Verdict chosen as slicing criterion. The closure slice (colored in green) identifies
all elements that may impact the execution of the
slicing criterion. In addition, all elements of the
other three sub-maps of Fig. 2 are also part of the
closure slice.
Later, Alkaf et al. (2019) proposed a forward
slicing technique to capture and grasp how changes
propagate within UCM models, and between GRL
and UCM models through URN traceability links.
The developed technique benefits maintainers
when locating the impact of potential modifications on both GRL and UCM model elements in
anticipation of carrying out a maintenance task
by increasing their productivity and cutting down
the cost of maintenance. Fig. 14 illustrates the
application of change impact analysis (CIA) to
the IMS example, with responsibility Store Verdict
chosen as slicing criterion. Potentially impacted
model elements are marked in green.
In addition, Fig. 15 illustrates the marked GRL
constructs (in purple color, as per the jUCMNav
CIA feature implementation) when we apply CIA
with task Fully Automated as slicing criterion. The
four impacted GRL strategies are also identified
(shown as URN comment). The Hospital actor is
not shown here as none of its elements are impacted. This support of GRL and UCM elements
as slicing criteria, with tool support for forward
and backward slicing, bring benefits to modellers
who need to understand complex process models
linked to goal models, and enables them to assess
the impact of model modifications before they are
implemented.
9.4 URN and Social Sciences Techniques
URN has also been integrated with techniques and
concepts from social sciences such as psychological paradigms and human values. Information
collected based on two psychological paradigms
and elicitation techniques, Activity Theory (Georg
et al. 2015) and Distributed Cognition (Hundal
and Mussbacher 2018), is systematically transformed into URN models for further analysis.
Both benefit from the combined modelling of
goals and processes in URN to fully describe and
consequently analyze the elicited findings, and the
explicit consideration of social science paradigms
helps discover requirements and process elements
that would be difficult to infer otherwise. The
Continual Value(s) Assessment (CVA) framework
(Perera et al. 2020) again uses goal and process
models but also feature models to systematically
integrate impact analysis on human values throughout software development.
10 Research Challenges and
Opportunities
This section discusses recent progress related to
generic URN-related challenges identified in the
literature a decade ago, and then discusses current
research challenges and opportunities related to
goal/process integration with URN.
10.1 A Decade of Improvements
In their literature review, Amyot and Mussbacher
(2011) highlighted eight research challenges for
the User Requirements Notation:
1. Domain-Specific Profiles: the ability to tailor
URN to specific domains was explored substantially in the past decade with different profiles in
the legal domain, as discussed in Sect. 5 with the
Legal-GRL profile (and its Variability-Aware
variant), the Indicator-based Policy Compliance
Framework, and the Goal-oriented Regulatory
Intelligence Method. Other domains where new
profiles were developed include healthcare with
the Activity-based Process Integration framework (Sect. 6) and value co-creation in service
systems (Sect. 7). Although tool support in
jUCMNav enables stereotypes for concepts and
relationships, well-formedness constraints, and
analysis algorithms to be tailored in profiles, the
symbols and terminology used are still based
on the URN standard, so part of the initial
challenge remains.
2. Enhanced Workflow Executions: The 2012
version of the URN standard included many
enhancements to the requirements for workflow
execution (called UCM path traversal mechanism), including new types of dynamic stubs
International Journal of Conceptual Modeling
Vol. 17, No. 2 (2022). DOI:10.18417/emisa.17.2
28 D. Amyot, O. Akhigbe, M. Baslyman, S. Ghanavati, M. Ghasemi, J. Hassine, L. Lessard, et al.
Special Issue on 100 Years of Graphical Business Process Modelling by R. Laue, H. Mayr & B. Thalheim
Figure 13: Closure slice of the IMS root map with respect to the slicing criterion SC = (Store Verdict, -).
Figure 14: Marked IMS root map with respect to the slicing criterion SC = (Store Verdict, -).
Figure 15: IMS: Marked GRL elements (extract).
Enterprise Modelling and Information Systems Architectures
Vol. 17, No. 2 (2022). DOI:10.18417/emisa.17.2
Combining Goal modelling with Business Process modelling 29
Special Issue on 100 Years of Graphical Business Process Modelling by R. Laue, H. Mayr & B. Thalheim
(with synchronization and blocking behaviour),
precise handling of map instances, new types
of timers (transient/persistent), and support for
exceptions. Many of these improvements made
their way to jUCMNav, but a proper debugging
environment is still missing. Non-standard execution mechanisms involving aspect-orientation
(Sect. 9.2) and slicing (Sect. 9.3) were also explored, with tool support.
3. Performance Management: On the GRL side,
indicators were added to the standard in 2012
to support performance management in the
goal view, with support in jUCMNav. Other
non-standard tool features include support for
qualitative KPIs, trend computation, and various export and reporting facilities (Amyot et al.
2012). This represented a major advancement
to goal-oriented modelling languages.
4. Compliance management: There are major
results in this area, as highlighted by the publication trends observed in Fig.7. Many such
results are discussed in Sect. 5 for legal and
regulatory compliance of business processes
and of organizations, but also in Sect. 4 in
terms of alignment between processes and organizational objectives, as well as in Sect. 6
for business and clinical process adaptation
contexts.
5. Formal Semantics: Although the URN standard does not include formal semantics, one
based on arithmetic functions was introduced
for GRL (Fan et al. 2018). Some new formal
semantics were also introduced for UCM beyond those that previously existed (covered in
Sect. 2.1), e. g., based on Object-Z (Dongmo
and Van der Poll 2014) or Maude (Mokhati
and Menassel 2013). On the UCM side, these
formal semantics still remain partial however
as they do not capture all the concepts and
informal workflow semantics of the standard.
6. Improved Analysis Techniques: There were major contributions in this area, especially in terms
of goal model analysis with support for indicators and constraint-based optimizations, support for enhanced workflow executions (point 2
above), analysis techniques exploiting domainspecific profiles (point 1 above), and other
combined UCM/GRL analysis techniques highlighted in Sections 4 to 9.
7. Improved Model Representations: Although
the graphical syntax of URN has not really
evolved in the past decade, one major addition
in the 2018 edition of the URN standard is
a textual syntax (see Sect. 2.3) enabling the
coding of models using conventional text-based
environments.
8. Guidelines and Methodologies: The URN standard focuses on the description of the language
rather than on methodological aspects. This
general lack of guidelines was addressed by
most of the frameworks discussed in Sections 5
to 8, which describe how to use URN or one of
its domain-specific profiles to support conceptual modelling and analysis activities.
As shown above, many generic URN-related
challenges were partially addressed by the community in the last decade. However, there are new
ones that have appeared since, especially around
business process/goal integration.
10.2 Remaining Challenges and
Opportunities
Although much progress has been made in enabling and supporting the combined used of process modelling and goal modelling with the User
Requirements Notation, many research challenges
and opportunities remain. Each of the areas covered in the preceding sections are ongoing, and
research challenges are detailed in their respective
referenced sources.
In general, while URN has seen many applications in academia and in industry, further work and
education is required to gain wider mainstream
adoption by practitioners. One approach would be
to further customize UCM and GRL to the needs
of specific domains (e. g., financial industry) so
that domain-specific terminology and syntax can
be used, and reusable assets can be developed.
This would go beyond the labeling-based profiles
International Journal of Conceptual Modeling
Vol. 17, No. 2 (2022). DOI:10.18417/emisa.17.2
30 D. Amyot, O. Akhigbe, M. Baslyman, S. Ghanavati, M. Ghasemi, J. Hassine, L. Lessard, et al.
Special Issue on 100 Years of Graphical Business Process Modelling by R. Laue, H. Mayr & B. Thalheim
developed so far, and require better tool support
than what is currently offered in jUCMNav.
Some challenges are shared with other goal
modelling languages. Scalability, i. e., the ability to handle large models, is a common issue in
goal modelling, and can potentially be addressed
by language design (e. g., use of aspect-oriented
constructs), in the usage methods, as well as by innovative tool support. Although jUCMNav allows
GRL models to include many views as separate
diagrams and to generate new diagrams from
model elements by expanding neighbouring elements, modularity of goal models is still an issue.
For URN, the loose coupling of two modelling
paradigms also presents some additional research
challenges and opportunities.
The primarily manual task of coordinating and
maintaining traceability between goals and processes, which is at the core of many of the above
techniques for alignment, compliance, adaptation,
and value co-creation, can be time-consuming
and error-prone. Automating some aspects of
these tasks (e. g., identifying goal/process relations through natural language processing (Aung
et al. 2020), and checking and maintaining semantic correctness and consistency) would boost
productivity and facilitate practical adoption.
Automation of many manual activities beyond
traceability is also an important research direction,
especially when usability and integration into
existing methods used by practitioners must be
taken into consideration.
On the methods side, there is potential synergy between URN modelling and performance
management methods, as the former can tap into
existing KPI definition and data collection processes, whereas the latter can benefit from the
analytical and reasoning power of URN to better
enable decision making and consensus building.
One significantly new research direction for
URN is to provide support for processes that
engage with fluid social networks and evolving
ecosystems, as well as with socio-cyber-physical
systems that include multi-channel heterogeneous
content from communities with shifting identities
and disparate value systems, as human and automated systems react to each other within ever
shortening time frames, leaving little time for
rational decision making.
Another significant research direction will be to
consider how to integrate the use of advanced machine learning algorithms (including deep learning) that draw on massive data sets and diverse
data streams (from the cloud, sensors, and social
and mobile networks) and could enhance the current business intelligence research in the areas
of regulatory compliance and business process
adaptation/improvement.
On a broader scope, as digitalization is rapidly
transforming every business sector and almost
all aspects of life, we expect that the past two
decades of experiences with URN will provide
the foundations for addressing many of the challenges that we face today. Processes will need
to be more dynamic and agile to adapt to rapid
and continual change. Associating and aligning
process design with goal reasoning, as supported
by URN, provides a foundational building block
for model-based transformation.
11 Conclusion
One century after the seminal work of Gilbreth
and Gilbreth (1921), process modelling is more
popular and diverse than ever. Goal modelling,
although much younger, has shown interesting
potential in complementing process modelling for
many important types of activities.
This paper gives a contemporary overview of
the User Requirements Notation standard, of its
general usage, and of the benefits brought by
its combination of GRL and UCM modelling.
Although many combinations of existing goal and
process modelling languages have been the topic
of research activities and practical applications,
few combinations have been as well studied in
diverse contexts as GRL combined with UCM.
Sect. 3 mapped the academic interest in URN
modelling and the many application areas where
URN was used over time.
Enterprise Modelling and Information Systems Architectures
Vol. 17, No. 2 (2022). DOI:10.18417/emisa.17.2
Combining Goal modelling with Business Process modelling 31
Special Issue on 100 Years of Graphical Business Process Modelling by R. Laue, H. Mayr & B. Thalheim
The paper answers the posed research question “What are the benefits of combining goal
modelling with process modelling?” by examining over two decades of contributions of URNbased techniques to the domains of goal/process
alignment, regulatory compliance and intelligence,
process adaptation and improvement, value cocreation and service systems, and goal-oriented
process mining. Advanced URN-based modelling
and analysis techniques for features, aspects, slicing, and concepts from social sciences are also
discussed.
As a summary of observed benefits, combined
goal and process modelling in URN helps in the
following ways:
• For the alignment of goals and processes
(Sect. 4), this combination helps determine
i) whether existing processes as they are defined meet organization and stakeholder goals
and ii) whether goals are properly covered by
such processes even when both goals and processes evolve.
• Regarding regulatory compliance and intelligence (Sect. 5), a goal/process combination
helps model both the intent and the structure
of laws and regulations as well as operational
aspects (high-level processes) i) to better ensure
that organization processes comply with laws,
ii) to maintain compliance when changes occur
to laws or organization processes and objectives,
iii) to allow experts to discuss ambiguities and
resolve them, iv) to detect hotspots in processes
regarding run-time regulatory compliance, and
v) to manage regulations in terms of continuous
monitoring, assessment, and reporting on their
efficiency and effectiveness.
• For process adaptation and improvement
(Sect. 6), we observe that goal/process modelling helps i) capture the dynamic influence of
one view on the other, ii) model feedback loops
of self-adaptive socio-technical systems based
on rules, run-time goal analysis, and a process
execution/simulation environment, iii) capture
and incrementally evaluate the impact of each
newly integrated activity on stakeholder goals
to select the best process integration alternative,
and iv) formalize social concerns and measures
during change management of processes.
• For value co-creation and service systems
(Sect. 7), this combination helps i) define and
measure a service’s value from varied stakeholder perspectives in relation to processes,
ii) redefine the endpoint of a service process
as the moment when it creates value for stakeholders, and iii) support the identification of
additional system requirements that allow organizational actors to successfully integrate a
service.
• For goal-oriented process mining (Sect. 8), combined goal and process modelling helps improve
the precision, rationality, and interpretability
of the process models discovered through process mining by extracting simpler and understandable process models that meet measurable
goal-oriented criteria and eventually improve
the satisfaction of process stakeholders.
• Regarding more advanced URN-based modelling techniques (Sect. 9):
– Support for features and aspects in combined
goal and process models improves i) the accuracy of goal model analysis on which the
selection of (potentially reusable) features
and processes is based and ii) the coordinated, semantics-based reuse across different
types of models.
– Support for slicing in combined goal/process modelling helps modellers who need to
understand complex process models linked
to goal models, and enables them to assess
the impact of model modifications before
implementation.
– This combination also supports the more
complete description and subsequent analysis
of concepts from social sciences (e. g., information resulting from Activity Theory and
Distributed Cognition, or based on human
values), which in turn promote the discovery
of requirements and process elements that
would be difficult to infer otherwise.
International Journal of Conceptual Modeling
Vol. 17, No. 2 (2022). DOI:10.18417/emisa.17.2
32 D. Amyot, O. Akhigbe, M. Baslyman, S. Ghanavati, M. Ghasemi, J. Hassine, L. Lessard, et al.
Special Issue on 100 Years of Graphical Business Process Modelling by R. Laue, H. Mayr & B. Thalheim
Although some of the above benefits might be
unique to URN, we suspect that many benefits
can be generalized to other combinations of goal
and process modelling languages. Modelling with
two complementary notations does not only bring
benefits. In addition to the extra effort required
to combine two such partial but complementary
views, important challenges and exciting opportunities for future research are discussed in Sect. 10,
focusing on the areas of adoption, usable automation, data-driven applications (e.g, AI/machine
learning), and socio-cyber-physical systems.
References
Abrahão S., Insfrán E., González-Ladrón-deGuevara F., Fernández-Diego M., Cano-Genoves
C., de Oliveira R. P. (2019) Assessing the effectiveness of goal-oriented modeling languages: A
family of experiments. In: Information and Software Technology 116, p. 106171
Akhigbe O., Alhaj M., Amyot D., Badreddin O.,
Braun E., Cartwright N., Richards G., Mussbacher
G. (2014) Creating Quantitative Goal Models:
Governmental Experience. In: Conceptual Modeling: 33rd International Conference, ER 2014.
Springer, pp. 466–473
Akhigbe O., Amyot D., Anda A. A., Lessard L.,
Xiao D. (2016) Consistency Analysis for User
Requirements Notation Models. In: Proceedings
of the Ninth International i* Workshop (iStar
2016). CEUR-WS Vol. 1674, pp. 43–48
Akhigbe O., Amyot D., Richards G. (2019) A
systematic literature mapping of goal and nongoal modelling methods for legal and regulatory
compliance. In: Requirements Engineering 24(4),
pp. 459–481
Akhigbe O., Heap S., Islam S., Amyot D., Mylopoulos J. (2017) Goal-Oriented Regulatory Intelligence: How Can Watson Analytics Help? In:
Conceptual Modeling (ER 2017). Lecture Notes
in Computer Science Vol. 10650. Springer, pp. 77–
91
Akhigbe O. S. (2018) A Goal-Oriented Method
for Regulatory Intelligence. PhD thesis, University
of Ottawa, Canada
Alkaf H. S., Hassine J., Binalialhag T., Amyot
D. (2019) An automated change impact analysis
approach for User Requirements Notation models.
In: Journal of Systems and Software 157
Amyot D., Becha H., Bræk R., Rossebø J. E. Y.
(2008) Next generation service engineering. In:
2008 First ITU-T Kaleidoscope Academic Conference. IEEE, pp. 195–202
Amyot D., Ghanavati S., Horkoff J., Mussbacher
G., Peyton L., Yu E. (2010) Evaluating goal models
within the goal-oriented requirement language. In:
International Journal of Intelligent Systems 25 (8),
pp. 841–877
Amyot D., Logrippo L. (2000) Use Case Maps
and Lotos for the prototyping and validation of
a mobile group call system. In: Computer Communications 23(12), pp. 1135–1157
Amyot D., Mussbacher G. (2011) User Requirements Notation: the first ten years, the next ten
years. In: Journal of Software (JSW) 6(5), pp. 747–
768
Amyot D., Shamsaei A., Kealey J., Tremblay E.,
Miga A., Mussbacher G., Alhaj M., Tawhid R.,
Braun E., Cartwright N. (2012) Towards Advanced
Goal Model Analysis with jUCMNav. In: Advances in Conceptual Modeling. Springer, pp. 201–
210
Amyot D., Yan J. B. (2008) Flexible Verification of
User-Defined Semantic Constraints in Modelling
Tools. In: Proceedings of CASCON 2008. ACM,
pp. 81–95
Anda A., Amyot D. (2020) An Optimization Modeling Method for Adaptive Systems Based on Goal
and Feature Models. In: Proceedings of the Tenth
International Model-Driven Requirements Engineering Workshop (MoDRE). IEEE, pp. 10–20
Androutsopoulos K., Clark D., Harman M., Krinke
J., Tratt L. (2013) State-Based Model Slicing: A
Survey. In: ACM Computing Surveys 45(4), pp. 1–
36
Enterprise Modelling and Information Systems Architectures
Vol. 17, No. 2 (2022). DOI:10.18417/emisa.17.2
Combining Goal modelling with Business Process modelling 33
Special Issue on 100 Years of Graphical Business Process Modelling by R. Laue, H. Mayr & B. Thalheim
Antón A. I. (1996) Goal-based requirements analysis. In: Proceedings of the Second International
Conference on Requirements Engineering. IEEE,
pp. 136–144
Aung T. W. W., Huo H., Sui Y. (2020) A Literature
Review of Automatic Traceability Links Recovery
for Software Change Impact Analysis. In: 28th
International Conference on Program Comprehension (ICPC 2020). ACM, pp. 14–24
Baslyman M., Almoaber B., Amyot D., Bouattane E. M. (2017a) Activity-based Process Integration in Healthcare with the User Requirements Notation. In: International Conference on
E-Technologies. Springer, pp. 151–169
Baslyman M., Almoaber B., Amyot D., Bouattane
E. M. (2017b) Using Goals and Indicators for
Activity-based Process Integration in Healthcare.
In: Procedia Computer Science 113, pp. 318–325
Baslyman M., Amyot D., Alshalahi Y. (2019)
Lean healthcare processes: effective technology
integration and comprehensive decision support
using requirements engineering methods. In: 2019
IEEE/ACM 1st International Workshop on Software Engineering for Healthcare (SEH). IEEE,
pp. 37–44
Behnam S. A., Amyot D. (2013) Evolution mechanisms for goal-driven pattern families used in
business process modelling. In: International Journal of Electronic Business 10(3), pp. 254–291
Binalialhag T., Hassine J., Amyot D. (2018) Static
slicing of Use Case Maps requirements models. In:
Software & Systems Modeling 18(4), pp. 2465–
2505
Bitner M. J., Ostrom A. L., Morgan F. N. (2008)
Service Blueprinting: A Practical Technique for
Service Innovation. In: California Management
Review 50(3), pp. 66–94
Bresciani P., Perini A., Giorgini P., Giunchiglia
F., Mylopoulos J. (2004) Tropos: An agentoriented software development methodology. In:
Autonomous Agents and Multi-Agent Systems
8(3), pp. 203–236
Buhr R. J. A., Casselman R. S. (1995) Use Case
Maps for Object-Oriented Systems. Prentice-Hall
Buhr R. J. A. (1998) Use case maps as architectural
entities for complex systems. In: IEEE Transactions on Software Engineering 24(12), pp. 1131–
1155
Cardoso E. C. S., Guizzardi R. S. S., Almeida
J. P. A. (2011) Aligning goal analysis and business
process modelling: a case study in healthcare. In:
International Journal of Business Process Integration and Management 5(2), pp. 144–158
Chung L., Nixon B. A., Yu E., Mylopoulos J.
(2000) Non-functional requirements in software
engineering Vol. 5. Springer
Clausing D., Hauser J. R. (1988) The house of quality. In: Harvard Business Review 66(3), pp. 63–
73
Dalpiaz F., Franch X., Horkoff J. (2016) iStar 2.0
Language Guide. In: CoRR
Decreus K., Snoeck M., Poels G. (2009) Practical
Challenges for Methods Transforming i* Goal
Models into Business Process Models. In: 2009
17th IEEE International Requirements Engineering Conference. IEEE, pp. 15–23
Dongmo C., Van der Poll J. A. (2014) Addressing
the construction of Z and Object-Z with use case
maps (UCMs). In: International Journal of Software Engineering and Knowledge Engineering
24(2), pp. 285–327
Dotan D., Rubin J., Yatzkar-Haham T. (2013)
Configuration management system for software
product line development environment. US Patent
8,549,473
Duran M. B. (2018) Reusable Goal Models. PhD
thesis, McGill University, Canada
Duran M. B., Mussbacher G. (2016) Investigation
of feature run-time conflicts on goal model-based
reuse. In: Information Systems Frontiers, pp. 1–21
Duran M. B., Mussbacher G. (2018) Top-Down
Evaluation of Reusable Goal Models. In: New Opportunities for Software Reuse. Springer, pp. 76–
92
International Journal of Conceptual Modeling
Vol. 17, No. 2 (2022). DOI:10.18417/emisa.17.2
34 D. Amyot, O. Akhigbe, M. Baslyman, S. Ghanavati, M. Ghasemi, J. Hassine, L. Lessard, et al.
Special Issue on 100 Years of Graphical Business Process Modelling by R. Laue, H. Mayr & B. Thalheim
Duran M. B., Mussbacher G. (2019) Reusability
in goal modeling: A systematic literature review.
In: Information and Software Technology 110,
pp. 156–173
Fan Y., Anda A. A., Amyot D. (2018) An Arithmetic Semantics for GRL Goal Models with Function Generation. In: International Conference on
System Analysis and Modeling. Springer, pp. 144–
162
Franchitti J.-C. L. (2019) Active adaptation of
networked compute devices using vetted reusable
software components. US Patent 10,338,913
Genon N., Amyot D., Heymans P. (2011)
Analysing the Cognitive Effectiveness of the
UCM Visual Notation. In: System Analysis and
Modeling: About Models (SAM 2010). Springer,
pp. 221–240
Georg G., Mussbacher G., Amyot D., Petriu D.,
Troup L., Lozano-Fuentes S., France R. (2015)
Synergy between Activity Theory and goal/scenario modeling for requirements elicitation, analysis, and evolution. In: Information and Software
Technology 59, pp. 109–135
Ghanavati S. (2013) Legal-URN Framework for
Legal Compliance of Business Processes. PhD
thesis, University of Ottawa, Canada
Ghanavati S., Amyot D., Rifaut A. (2014a) Legal Goal-Oriented Requirement Language (Legal
GRL) for Modeling Regulations. In: Proceedings
of the 6th International Workshop on Modeling in
Software Engineering. MiSE 2014. ACM, pp. 1–6
Ghanavati S., Rifaut A., Dubois E., Amyot D.
(2014b) Goal-oriented compliance with multiple
regulations. In: 2014 IEEE 22nd International Requirements Engineering Conference (RE). IEEE,
pp. 73–82
Ghasemi M., Amyot D. (2019a) Data Preprocessing for Goal-Oriented Process Discovery. In: 27th
IEEE International Requirements Engineering
Conference Workshops (RE 2019). IEEE, pp. 200–
206
Ghasemi M., Amyot D. (2019b) Goal-oriented
Process Enhancement and Discovery. In: International conference on business process management. Springer, pp. 102–118
Ghasemi M., Amyot D. (2020) From event logs
to goals: a systematic literature review of goaloriented process mining. In: Requirements Engineering 25(1), pp. 67–93
Gilbreth F. B., Gilbreth L. M. (1921) Process
charts. In: Annual Meeting of The American Society of Mechanical Engineers. New York
Godart C., Molli P., Canals G., Lonchamp J. (1992)
Extending cooperating transaction model to support software engineering activities. In: Proceedings of the Second International Conference on
Systems Integration, pp. 184–192
Gol Mohammadi N. (2019) Systematic Refinement of Trustworthiness Requirements using Goal
and Business Process Models In: Trustworthy
Cyber-Physical Systems Springer, pp. 119–144
Gonçalves E., Castro J., Araújo J., Heineck T.
(2018) A Systematic Literature Review of iStar
extensions. In: Journal of Systems and Software
137, pp. 1–33
Guizzardi R., Reis A. N. (2015) A Method to Align
Goals and Business Processes. In: Conceptual
Modeling. Springer, pp. 79–93
Hashmi M., Governatori G., Lam H., Wynn M. T.
(2018) Are we done with business process compliance: state of the art and challenges ahead.
In: Knowledge and Information Systems 57(1),
pp. 79–133
Hassine J., Amyot D. (2016) A questionnairebased survey methodology for systematically validating goal-oriented models. In: Requirements
Engineering 21(2), pp. 285–308
Hassine J., Amyot D. (2017) An empirical approach toward the resolution of conflicts in goaloriented models. In: Software & System Modeling
16(1), pp. 279–306
Enterprise Modelling and Information Systems Architectures
Vol. 17, No. 2 (2022). DOI:10.18417/emisa.17.2
Combining Goal modelling with Business Process modelling 35
Special Issue on 100 Years of Graphical Business Process Modelling by R. Laue, H. Mayr & B. Thalheim
Hassine J., Rilling J., Dssouli R. (2005a) Abstract
Operational Semantics for Use Case Maps. In:
Formal Techniques for Networked and Distributed
Systems - FORTE 2005, 25th IFIP WG 6.1 International Conference. Lecture Notes in Computer
Science Vol. 3731. Springer, pp. 366–380
Hassine J., Rilling J., Dssouli R. (2005b) An ASM
Operational Semantics for Use Case Maps. In:
13th IEEE International Conference on Requirements Engineering (RE 2005). IEEE, pp. 467–
468
Horkoff J., Aydemir F. B., Cardoso E., Li T., Maté
A., Paja E., Salnitri M., Piras L., Mylopoulos J.,
Giorgini P. (2019) Goal-oriented requirements engineering: an extended systematic mapping study.
In: Requirements Engineering 24(2), pp. 133–160
Horkoff J., Barone D., Jiang L., Yu E. S. K., Amyot
D., Borgida A., Mylopoulos J. (2014) Strategic
business modeling: representation and reasoning.
In: Software & Systems Modeling 13(3), pp. 1015–
1041
Hundal K. S., Mussbacher G. (2018) Model-Based
Development with Distributed Cognition. In: 2018
IEEE 8th International Model-Driven Requirements Engineering Workshop (MoDRE). IEEE,
pp. 26–35
IBM (2019) IBM ILOG CPLEX Optimization
Studio. https://ibm.co/2YRj8oy
IEEE (2012) ISO/IEC/IEEE 31320-1:2012 – Information technology – Modeling Languages–Part 1:
Syntax and Semantics for IDEF0. https://www.iso.
org/standard/60615.html
Insfrán E., Abrahão S., de Oliveira R. P., GonzálezLadrón-de-Guevara F., Fernández-Diego M.,
Cano-Genoves C. (2017) Specifying Value in GRL
for Guiding BPMN Activities Prioritization. In:
Information Systems Development (ISD 2017)
Intersoft Consulting (2016) General Data Protection Regulation (GDPR). https://gdpr-info.eu/
ISO (2019) ISO/IEC 15909-1:2019 – Systems and
software engineering – High-level Petri nets – Part
1: Concepts, definitions and graphical notation.
https://www.iso.org/obp/ui/#iso:std:iso-iec:15909:-
1:en
ITU-T (2011) Recommendation Z.150 (02/11)
User Requirements Notation (URN) - Language
requirements and framework. http://www.itu.int/rec/
T-REC-Z.150/en
ITU-T (2018) Recommendation Z.151 (10/18)
User Requirements Notation (URN) – Language
definition. http://www.itu.int/rec/T-REC-Z.151/en
Jacobs S., Holten R. (1995) Goal Driven Business
Modelling: Supporting Decision Making within
Information Systems Development. In: ACM,
pp. 96–105
Karagiannis D., Mayr H. C., Mylopoulos J. (eds.)
Domain-Specific Conceptual Modeling, Concepts,
Methods and Tools. Springer
Koliadis G., Ghose A. (2006) Relating Business Process Models to Goal-Oriented Requirements Models in KAOS. In: Advances in Knowledge Acquisition and Management (PKAW 2006).
Springer, pp. 25–39
Kumar R., Mussbacher G. (2018) Textual User
Requirements Notation. In: System Analysis and
Modeling (SAM 2018). Springer, pp. 163–182
Kuziemsky C., Liu X., Peyton L. (2010) Leveraging goal models and performance indicators to
assess health care information systems. In: Quality
of Information and Communications Technology
(QUATIC), 2010 Seventh International Conference on the. IEEE, pp. 222–227
Lessard L. (2015a) An Intentional Approach to
the Engineering of Knowledge-intensive Service
Systems. In: Procedia Manufacturing 3 AHFE
2015, pp. 3383–3390
Lessard L. (2015b) Modeling Value Cocreation
Processes and Outcomes in Knowledge-Intensive
Business Services Engagements. In: Service Science 7(3), pp. 181–195
International Journal of Conceptual Modeling
Vol. 17, No. 2 (2022). DOI:10.18417/emisa.17.2
36 D. Amyot, O. Akhigbe, M. Baslyman, S. Ghanavati, M. Ghasemi, J. Hassine, L. Lessard, et al.
Special Issue on 100 Years of Graphical Business Process Modelling by R. Laue, H. Mayr & B. Thalheim
Lessard L., Amyot D., Aswad O., Mouttham A.
(2020) Expanding the nature and scope of requirements for service systems through ServiceDominant Logic: the case of a telemonitoring service. In: Requirements Engineering 25(3), pp. 273–
293
Liaskos S., Jalman R., Aranda J. (2012) On eliciting contribution measures in goal models. In:
2012 20th IEEE International Requirements Engineering Conference (RE). IEEE, pp. 221–230
Liscano R., Fullarton S., Mussbacher G., Gray T.
(2008) System and method for remote assembly of
messages to create a control message. US Patent
7,318,109
Liu L., Yu E. (2004) Designing information systems in social context: a goal and scenario modelling approach. In: Information Systems 29(2),
pp. 187–203
Liu Y., Su Y., Yin X., Mussbacher G. (2014)
Combined propagation-based reasoning with goal
and feature models. In: 2014 IEEE 4th International Model-Driven Requirements Engineering
Workshop (MoDRE). IEEE, p. 2736
Lo A., Yu E. S. K. (2007) From Business Models
to Service-Oriented Design: A Reference Catalog
Approach. In: Conceptual Modeling (ER 2007).
Lecture Notes in Computer Science Vol. 4801.
Springer, pp. 87–101
Massey A., Holtgrefe E., Ghanavati S. (2017) Modeling Regulatory Ambiguities for Requirements
Analysis. In: Conceptual Modeling (ER2017,
pp. 231–238
Massey A., Rutledge R., Antón A., Swire P. (2014)
Identifying and Classifying Ambiguity for Regulatory Requirements. In: 2014 IEEE 22nd International Requirements Engineering Conference
(RE). IEEE, pp. 83–92
Mears M. G., Nierzwick M., Pash P. E., Rizzo
V. R., Steiner B., Westerfield K. M. (2016) Handheld diabetes manager with a user interface for
displaying a status of an external medical device.
US Patent 9,252,870
Mili H., Tremblay G., Jaoude G. B., Lefebvre
É., Elabed L., Boussaidi G. E. (2010) Business
Process Modeling Languages: Sorting through
the Alphabet Soup. In: ACM Computing Surveys
43(1)
Mokhati F., Menassel Y. (2013) Towards formalising use case maps in Maude strategy language: application to multi-agent systems. In: International
journal of computer applications in technology
47(2-3), pp. 138–151
Moody D. L., Heymans P., Matulevicius R. (2010)
Visual syntax does matter: improving the cognitive effectiveness of the i* visual notation. In:
Requirements Engineering 15(2), pp. 141–175
Mussbacher G., Amyot D. (2008) Assessing the
Applicability of Use Case Maps for Business Process and Workflow Description. In: 2008 International MCETECH Conference on e-Technologies
(MCETECH 2008). IEEE, pp. 219–222
Mussbacher G. (2010) Aspect-oriented User Requirements Notation. PhD thesis, University of
Ottawa, Canada
Mussbacher G., Amyot D., Araújo J., Moreira A.
(2010a) Requirements Modeling with the Aspectoriented User Requirements Notation (AoURN): A
Case Study In: Transactions on Aspect-Oriented
Software Development VII: A Common Case
Study for Aspect-Oriented Modeling Springer,
pp. 23–68
Mussbacher G., Amyot D., Whittle J. (2009)
Refactoring-Safe Modeling of Aspect-Oriented
Scenarios. In: Model Driven Engineering Languages and Systems. Springer, pp. 286–300
Mussbacher G., Amyot D., Whittle J. (2013) Composing Goal and Scenario Models with the AspectOriented User Requirements Notation Based on
Syntax and Semantics In: Aspect-Oriented Requirements Engineering Springer, pp. 77–99
Mussbacher G., Araújo J., Moreira A., Amyot
D. (2012) AoURN-based modeling and analysis
of software product lines. In: Software Quality
Journal 20(3–4), pp. 645–687 10.1007/s11219-
011-9153-8
Enterprise Modelling and Information Systems Architectures
Vol. 17, No. 2 (2022). DOI:10.18417/emisa.17.2
Combining Goal modelling with Business Process modelling 37
Special Issue on 100 Years of Graphical Business Process Modelling by R. Laue, H. Mayr & B. Thalheim
Mussbacher G., Whittle J., Amyot D. (2010b)
Modeling and detecting semantic-based interactions in aspect-oriented scenarios. In: Requirements Engineering Journal (REJ) 15(2), pp. 197–
214
Mylopoulos J., Borgida A., Jarke M., Koubarakis
M. (1990) Telos: Representing Knowledge about
Information Systems. In: ACM Transactions on
Information Systems 8(4), pp. 325–362
Mylopoulos J., Chung L., Nixon B. A. (1992) Representing and using nonfunctional requirements: a
process-oriented approach. In: IEEE Transactions
on Software Engineering 18(6), pp. 483–497
Ni Q., Bertino E., Lobo J., Brodie C., Karat C.,
Karat J., Trombetta A. (2010) Privacy-aware rolebased access control. In: ACM Transactions on
Information and System Security 13(3), 24:1–
24:31
Odeh Y., Green S., Odeh M. (2018) Deriving
Goal-Oriented Models from Business Process
Models: Applied to Cancer Care Organisation. In:
2018 1st International Conference on Cancer Care
Informatics (CCI). IEEE, pp. 125–135
OMG (2014) Business Process Model and Notation (BPMN), version 2.0.2. https://www.omg.org/
spec/BPMN/2.0.2
OMG (2017) Unified Modeling Language (UML),
version 2.5.1. https://www.omg.org/spec/UML/2.5.1
Patrício L., de Pinho N. F., Teixeira J. G., Fisk
R. P. (2018) Service Design for Value Networks:
Enabling Value Cocreation Interactions in Healthcare. In: Service Science 10(1), pp. 76–97
Perera H., Mussbacher G., Hussain W., Ara Shams
R., Nurwidyantoro A., Whittle J. (2020) Continual
Human Value Analysis in Software Development:
A Goal Model Based Approach. In: 2020 IEEE
28th International Requirements Engineering Conference (RE), pp. 192–203
Pesic M., Van der Aalst W. M. (2006) A declarative
approach for flexible business processes management. In: International conference on business
process management. Springer, pp. 169–180
Petriu D. B., Amyot D., Woodside C. M. (2003)
Scenario-Based Performance Engineering with
UCMNAV. In: System Design, 11th International
SDL Forum (SDL 2003). Lecture Notes in Computer Science Vol. 2708. Springer, pp. 18–35
Pinciroli F., Barros Justo J. L., Forradellas R.
(2020) Systematic mapping study: On the coverage
of aspect-oriented methodologies for the early
phases of the software development life cycle. In:
Journal of King Saud University - Computer and
Information Sciences
Rahman A., Amyot D. (2014) A DSL for importing models in a requirements management
system. In: IEEE 4th International Model-Driven
Requirements Engineering Workshop (MoDRE
2014). IEEE, pp. 37–46
Roy J.-F., Kealey J., Amyot D. (2006) Towards Integrated Tool Support for the User Requirements
Notation. In: System Analysis and Modeling: Language Profiles. Springer, pp. 198–215
Sartoli S., Ghanavati S., Siami Namin A. (2020a)
Compliance Requirements Checking in Variable
Environments. In: 2020 IEEE 44th Annual Computers, Software, and Applications Conference
(COMPSAC). IEEE, pp. 1093–1094
Sartoli S., Ghanavati S., Siami Namin A. (2020b)
Towards Variability-Aware Legal-GRL Framework for Modeling Compliance Requirements. In:
2020 IEEE 7th International Workshop on Evolving Security Privacy Requirements Engineering
(ESPRE). IEEE, pp. 7–12
Shamsaei A. (2012) Indicator-based policy compliance of business processes. PhD thesis, University
of Ottawa, Canada
Shamsaei A., Pourshahid A., Amyot D. (2010)
Business Process Compliance Tracking Using
Key Performance Indicators. In: Business Process
Management Workshops (BPM 2010). Lecture
Notes in Business Information Processing Vol. 66.
Springer, pp. 73–84
International Journal of Conceptual Modeling
Vol. 17, No. 2 (2022). DOI:10.18417/emisa.17.2
38 D. Amyot, O. Akhigbe, M. Baslyman, S. Ghanavati, M. Ghasemi, J. Hassine, L. Lessard, et al.
Special Issue on 100 Years of Graphical Business Process Modelling by R. Laue, H. Mayr & B. Thalheim
Sousa H. P., Leite J. C. S. d. P. (2014) Modeling
Organizational Alignment. In: Conceptual Modeling - 33rd International Conference (ER2014).
Lecture Notes in Computer Science Vol. 8824.
Springer, pp. 407–414
U.S. Government (1996) Health Insurance Portability and Accountability Act (HIPAA). http://www.
hhs.gov/ocr/privacy/
Van Der Aalst W. (2016) Data science in action.
In: Process mining. Springer, pp. 3–23
van Lamsweerde A. (2001) Goal-oriented requirements engineering: a guided tour. In: Proceedings
Fifth IEEE International Symposium on Requirements Engineering. IEEE, pp. 249–262
van Lamsweerde A. (2009) Requirements engineering: From system goals to UML models to
software Vol. 10. Chichester, UK: John Wiley &
Sons
Vigder M. (1992) Applying formal techniques
to the design of concurrent systems.. PhD thesis,
Carleton University, Canada
Vinay S., Aithal S., Sudhakara G. (2014) Effect
of Contribution Links on Choosing Hard Goals
in GORE Using AHP and TOPSIS. In: Emerging
Research in Electronics, Computer Science and
Technology: Proceedings of International Conference, ICERECT 2012. Springer, pp. 737–745
Vrbaski M., Mussbacher G., Petriu D. C., Amyot
D. (2012) Goal models as run-time entities in
context-aware systems. In: 7th Workshop on Models@run.time. ACM, pp. 3–8
Weigert T., Kolchin A., Potiyenko S., Gurenko
O., van den Berg A., Banas V., Chetvertak R.,
Yagodka R., Volkov V. (2019) Generating Test
Suites to Validate Legacy Systems. In: System
Analysis and Modeling, 11th International Conference (SAM 2019). Lecture Notes in Computer
Science Vol. 11753. Springer, pp. 3–23
Weigert T. J. (2017) System and method for iterative generating and testing of application code.
US Patent 9,778,916
Yu E., Amyot D., Mussbacher G., Franch X.,
Castro J. (2013) Practical applications of i* in
industry: The state of the art. In: Requirements
Engineering Conference (RE), 2013 21st IEEE
International. IEEE, pp. 366–367
Yu E. S. K. (1997) Towards modelling and reasoning support for early-phase requirements engineering. In: 3rd IEEE International Symposium
on Requirements Engineering (ISRE’97). IEEE,
pp. 226–235
Yu E. S., Giorgini P., Maiden N., Mylopoulos J.
(2011) Social modeling for requirements engineering. MIT press